\documentclass[11pt,a4paper]{article}
\usepackage[T1]{fontenc}
\usepackage[utf8]{inputenc}
\usepackage{lmodern}
\usepackage[margin=26mm,headheight=14pt]{geometry}
\usepackage{amsmath,amssymb,amsthm,mathtools}
\usepackage{booktabs,longtable,array}
\usepackage{microtype}
\usepackage{xcolor}
\usepackage{fancyhdr}
\usepackage{enumitem}
\usepackage[hyphens]{url}
\usepackage{xurl}
\usepackage[colorlinks=true,linkcolor=blue!45!black,citecolor=blue!45!black,urlcolor=blue!45!black]{hyperref}
\usepackage{bookmark}
\hypersetup{pdftitle={HODGE_SIGN: Sign-only purity and the Hodge-index comparison},pdfauthor={HODGE_SIGN research run, 21 November 2026}}
\setlength{\parindent}{0pt}
\setlength{\parskip}{5pt plus 1pt minus 1pt}
\setlength{\emergencystretch}{2em}
\setlist{topsep=4pt,itemsep=2pt}
\pagestyle{fancy}
\fancyhf{}
\fancyhead[L]{\small HODGE\_SIGN}
\fancyhead[R]{\small Checked mathematical report}
\fancyfoot[C]{\thepage}
\newtheorem{theorem}{Theorem}[section]
\newtheorem{lemma}[theorem]{Lemma}
\newtheorem{proposition}[theorem]{Proposition}
\newtheorem{corollary}[theorem]{Corollary}
\theoremstyle{remark}
\newtheorem{remark}[theorem]{Remark}
\DeclareMathOperator{\Tr}{Tr}
\DeclareMathOperator{\CH}{CH}
\DeclareMathOperator{\NS}{NS}
\DeclareMathOperator{\Spec}{Spec}
\DeclareMathOperator{\End}{End}
\DeclareMathOperator{\Hom}{Hom}
\title{\textbf{HODGE\_SIGN}\\[5pt]
\large Sign-only purity, evanescent cohomology,\\
and the limits of the Hodge-index comparison}
\author{Research run HODGE\_SIGN--20261121\\
\small GPT-6 Astra Pro, root and three independent task agents}
\date{21 November 2026}
\begin{document}
\maketitle
\begin{abstract}
We reconstruct the sign-only mechanism as a complete theorem: a generic
Fourier-compatible hypersurface of an affine power has a Gysin remainder
whose signed virtual class is effective and contains a tensor power of the
lowest compactly supported cohomology.  Parseval and tensor amplification
give the expected weight bound; stratification and duality give the usual
finite-field absolute values.  The sign is cohomological parity, not
positive definiteness of an intersection form.  We compute the split
quadric and CM elliptic-power cases, track every supplied SA06 C.1 row,
give a counterexample to the displayed higher-cycle height formula in
cote 44, prove corrected versions of the openness remark in cote 34, and
separate the actual theta-lift divisibility bounds from integral Lefschetz
and Wilson matrices.  The latter comparison includes the normalized
universal ring, exact small-power Smith forms, and explicit lattice
counterexamples.  No general standard conjecture or unspecified higher
arithmetic index theorem is claimed.
\end{abstract}
\tableofcontents
\newpage
\part{Checked findings and explicitly cited background}


\section{The question decided, and the scope of the decision}

The sign-only method does give the requested finite-field absolute values.
Its sign is an effective multiplicity sign in a Grothendieck group of
Frobenius modules.  It is not the signature of an intersection pairing.
The missing geometric detail in the short description is supplied by
\emph{generic weak Lefschetz}, followed by cancellation through Gysin maps.
For a smooth affine $U$ of dimension $d$, the resulting effective module
contains $H_c^d(U,\mathcal F)^{\otimes k}$ and an actual Gysin kernel.
Finite Noether normalization makes its defining hypersurface compatible
with the path Fourier estimate.  This proves the middle-degree upper bound;
hyperplane induction, localization, and duality give the other degrees.

This reconstructs a proof of known weight statements, rather than claiming
that the Weil conjectures are new.  It answers the proposed identification
of methods: it uses the restriction/Gysin and evanescence machinery of
cote 16, scans 34--40, without constructing the positive involution or
establishing the cycle conjectures discussed in scans 49--50.  No new general
Hodge standard conjecture follows.  The distinction remains even in cases
where an index theorem is independently known.

The arithmetic proposed continuation fails in its unqualified form.
Norms and slopes do not choose the sign of a bilinear form; moreover,
the displayed higher-codimension formula in cote 44, scan 4, Theorem 5,
has an explicit counterexample under all its displayed hypotheses.
The valid consequence of divisibility plus the inspected slope inequality
is instead an inequality for the size of the divisibility, stated below.
For cote 34, the literal ambient group in the struck Remark 1.17 makes the
claim false; natural precise corrections are proved, including its soluble
nonabelian version.  For the final question, rational $\mathrm{SL}_2$ data
do not specify an integral lattice.  The ordinary contraction operator,
normalized Lefschetz lowering, Wilson matrices, and actual step-10 period
bounds are computed and distinguished.

\subsection*{Status convention and sources}
All propositions in Part I have complete proofs below, or are explicitly
identified as cited background.  Their check events identify author-side
checks, independent re-proofs, source inspections, and computations separately.
Part II records propositions not established here and the exact remaining
comparison data; they are not silently promoted to findings.  In particular,
an unspecified corrected Arthur/Lefschetz integral comparison is not refuted
by refuting its unnormalized version.

The stipulated Alp\"oge RH and Dirichlet-GRH statements are accepted as
theorems, exactly as the task requires.  No authenticity or build verdict is
made.  The global RH input is used for the shifted Tate factor
$\zeta(s-3)$; it is not used in the finite-field proof.  No numbered Masullo
statement is in the supplied material or used in any proof here.
The original archive \emph{images} are not supplied: the mathematical
corrections in \texttt{TRANSCRIPTION\_NOTES.tsv} are candidates for a second
reading, not claims to have deciphered handwriting.

\section{A precise sign-only theorem}
\label{sec:sign}

\subsection{Conventions and weight-free inputs}
Fix a prime power $q$, a prime $\ell\nmid q$, and an embedding
$\iota:\overline{\mathbf Q}_\ell\hookrightarrow\mathbf C$ for the
absolute values in question.  Cohomology means geometric \emph{compactly
supported} \(\ell\)-adic cohomology when the subscript $c$ is written.
$F$ denotes geometric Frobenius, so $F$ acts by $q$ on
$\overline{\mathbf Q}_\ell(-1)$.  For a rational point over
$\mathbf F_{q^n}$, its local Frobenius means Frobenius for that residue
field; if its underlying closed point has degree $a\mid n$, this is the
$(n/a)$th power of the degree-$a$ local Frobenius.

The inputs, all independent of a weight theorem, are the trace formula,
K\"unneth, affine Artin vanishing, Poincar\'e duality, and generic weak
Lefschetz.  Relevant foundational locators are SGA 4, XIV, Th\'eor\`eme 3.1
(affine vanishing), XVII, \S5.4.3 (K\"unneth), and XVIII (duality);
SGA $4\frac12$, \emph{Th\'eor\`emes de finitude}, 1.9 and 4.3 (generic
base change and duality); and the trace-formula normalization in
\cite[(1.5.1), (1.5.4), (2.3), (2.14)]{WeilI}.
The exact generic restriction theorem used is
\cite[Theorem 3.4, Corollary 3.4.1, pp.~180--181;
Appendix A.5, pp.~220--221]{KatzACT}.
Its Appendix proves the needed assertion using generic base change and
Artin vanishing; no purity assertion from the weight-theoretic parts of
that article is used.  This last separation is essential to avoiding a
circular proof.

\begin{lemma}[HS-01: effective power sums]
\label{lem:power}
Let $\lambda_1,\ldots,\lambda_s$ be distinct nonzero complex numbers and
$c_j$ positive integers.  Suppose $B>0$ and
\[
 \left|\sum_{j=1}^s c_j\lambda_j^n\right|
 \le C\exp(o(n)) B^n \qquad(n\ge1).
\]
Then $|\lambda_j|\le B$ for every $j$.
\end{lemma}
\begin{proof}
The generating power series of the sums is holomorphic on $|z|<B^{-1}$.
In a neighborhood of zero it is the rational function
\[
 \sum_j \frac{c_j\lambda_j z}{1-\lambda_j z}.
\]
At $z=\lambda_j^{-1}$ its $j$th term has a nonzero simple pole, and every
other term is regular.  Analytic continuation therefore precludes such a
pole inside the disk of convergence.  This gives the assertion.  After
coincident eigenvalues are combined, positive multiplicities prevent their
coefficients from cancelling.  The values of the power sums themselves
need not be positive.  Traces of powers depend on eigenvalues with algebraic
multiplicity, so no semisimplicity is required.
\end{proof}

\begin{theorem}[HS-02: tensor amplification]
\label{thm:tensor}
Let $V$ be a finite-dimensional Frobenius module and $a\in\mathbf R$.
For an unbounded set of positive integers $k$, suppose that an actual
Frobenius module $K_k$, a sign $e_k$, and numbers $R_{k,n}$ satisfy
\[
 e_kR_{k,n}=\operatorname{Tr}(F^n,V^{\otimes k})
                  +\operatorname{Tr}(F^n,K_k),\qquad n\ge1,
\]
and
\[
 |\iota R_{k,n}|\le C_k\exp(o_k(n))q^{n(ka+c_k)},
 \qquad c_k/k\longrightarrow0.
\]
Then every eigenvalue $\alpha$ of $F$ on $V$ satisfies
$|\iota\alpha|\le q^a$.
If the spectrum is stable under $\alpha\mapsto q^{2a}/\alpha$, it is pure
of modulus $q^a$.
\end{theorem}
\begin{proof}
For fixed $k$, the signed remainder is an effective power sum.  In
particular, $\alpha^k$ is one of its roots, with positive multiplicity.
Lemma~\ref{lem:power} gives
$|\iota\alpha|^k\le q^{ka+c_k}$.  Take $k$th roots and then let $k$ tend
to infinity.  Apply the same upper bound to $q^{2a}/\alpha$ for the final
assertion.  The limit in the extension degree $n$ is taken first, for each
fixed $k$; the constants $C_k$, dimensions of $K_k$, and later finite base
extensions may all depend on $k$.
\end{proof}

\begin{lemma}[HS-03: path Parseval estimate]
\label{lem:path}
Let $E$ be a finite field of cardinality $Q$, $A=E^d$, and $\psi$ a
nontrivial additive character.  For $k\ge2$, let $b_i:A\times A\to E$
be nondegenerate bilinear forms and let $f_i:A\to E$ be arbitrary
functions.  Put
\[
 P(x_1,\ldots,x_k)=\sum_{i=1}^{k-1}b_i(x_i,x_{i+1})
                         +\sum_{i=1}^k f_i(x_i)+c.
\]
If $|t_i(x)|\le B$, then, for every $h\ne0$,
\[
 \left|\sum_{A^k}\prod_i t_i(x_i)\psi(hP(x))\right|
       \le B^k Q^{d(k+1)/2}.
\]
The weights may depend on $h$, provided the same bound holds.  When all
$t_i=t$, character orthogonality gives
\[
 QN=M^k+R,\quad M=\sum_A t(x),\quad
 N=\sum_{P(x)=0}\prod_i t(x_i),\quad
 |R|\le B^k Q^{d(k+1)/2+1}.
\]
\end{lemma}
\begin{proof}
The matrix $K_{i,h}(x,y)=\psi(hb_i(x,y))$ satisfies
$K_{i,h}K_{i,h}^*=Q^d I$: summing the character of
$hb_i(x-x',y)$ over $y$ gives zero unless $x=x'$.  Its operator norm on
counting-measure $\ell^2(A)$ is $Q^{d/2}$.  Absorb the vertex phases into
the weights.  The path sum is a product of $k-1$ such matrices,
$k-2$ diagonal multipliers of norm at most $B$, and two endpoint vectors
of norms at most $BQ^{d/2}$.  Cauchy--Schwarz and the operator norm
inequality give $B^kQ^{d(k+1)/2}$.  A bilinear endpoint contraction instead
of a Hermitian one is handled by conjugating the first endpoint vector.
Finally
$\sum_{h\in E}\psi(hP)=Q\,1_{P=0}$; the zero mode is $M^k$, and the
$Q-1$ other modes give the claimed remainder bound.
\end{proof}

\subsection{The exact geometric sign}

\begin{proposition}[SG-08: effective Gysin remainder]
\label{prop:gysin}
Let $U/\mathbf F_q$ be smooth affine of pure dimension $d\ge1$, and
$\mathcal F$ a lisse $\overline{\mathbf Q}_\ell$-sheaf.  Put
$X=U^k$, $D=kd$, and $\mathcal F_k=\mathcal F^{\boxtimes k}$.
Suppose $i:Y\hookrightarrow X$ is a smooth affine hypersurface for which
\[
 H^a(X,\mathcal F_k^\vee)\longrightarrow
 H^a(Y,i^*\mathcal F_k^\vee)
\]
is an isomorphism for $a<D-1$ and injective for $a=D-1$.
Let $M_n$ be the trace sum of $\mathcal F$ on $U(\mathbf F_{q^n})$,
$N_{k,n}$ the trace sum of $i^*\mathcal F_k$ on $Y(\mathbf F_{q^n})$,
and $R_{k,n}=q^nN_{k,n}-M_n^k$.  Then
\begin{equation}\label{eq:effective}
 (-1)^{D-1}R_{k,n}
   =\operatorname{Tr}(F^n,H_c^d(U,\mathcal F)^{\otimes k})
      +\operatorname{Tr}(F^n,K_k),
\end{equation}
where the actual module is
\begin{equation}\label{eq:kernel}
 K_k=\ker\bigl(H_c^{D-1}(Y,i^*\mathcal F_k)(-1)
                    \longrightarrow H_c^{D+1}(X,\mathcal F_k)\bigr).
\end{equation}
\end{proposition}
\begin{proof}
Affine Artin vanishing for the dual lisse sheaf and Poincar\'e duality
give $H_c^j(U,\mathcal F)=0$ for $j<d$, and similarly the groups on
$X$ and $Y$ vanish below $D$ and $D-1$, respectively.
Dualizing the restriction in the hypothesis gives the Frobenius-equivariant
Gysin maps
\[
 H_c^j(Y,i^*\mathcal F_k)(-1)\longrightarrow
 H_c^{j+2}(X,\mathcal F_k),
\]
isomorphic for $j>D-1$ and surjective for $j=D-1$.
For clarity about the twist, the target is equivalently
$H_c^{j+2}(X,\mathcal F_k)(1)$ when the source is not twisted.

Write $T_b=H_c^{D+b}(X,\mathcal F_k)$ and
$W_a=H_c^{D-1+a}(Y,i^*\mathcal F_k)$, $a,b\ge0$.
The trace formula and $(-1)$ twist give
\[
 (-1)^{D-1}R_{k,n}
 =\sum_{a\ge0}(-1)^a\operatorname{Tr}(F^n,W_a(-1))
   +\sum_{b\ge0}(-1)^b\operatorname{Tr}(F^n,T_b).
\]
For $a\ge1$, Gysin identifies $W_a(-1)$ with $T_{a+1}$, whose sign in
the second sum is opposite.  The remaining virtual class is
$[W_0(-1)]+[T_0]-[T_1]$.  Surjectivity at the border turns this into
$[K_k]+[T_0]$.  In total K\"unneth degree $D=kd$, each of the $k$ factors
must have its minimum possible degree $d$, hence
$T_0=H_c^d(U,\mathcal F)^{\otimes k}$.
This is an equality of Frobenius-module classes, not merely dimensions.
\end{proof}

No concentration of the entire cohomology of $U$ or $Y$ in one degree was
assumed.  Higher groups cancel by actual Gysin isomorphisms; the border
surjection leaves an actual kernel.  In perverse normalization
$P=\mathcal F[d]$, form $M,N,R$ using the alternating stalk traces of
$P$ instead of those of $\mathcal F$.  The identity then has the single sign $-R$:
if $P_Y=i^*(P^{\boxtimes k})[-1]$, it is the trace of
\[
 H_c^0(U,P)^{\otimes k}\ \oplus\
 \ker\bigl(H_c^0(Y,P_Y)(-1)\longrightarrow
                         H_c^1(X,P^{\boxtimes k})\bigr).
\]
The minus sign comes from $i^*(P^{\boxtimes k})=P_Y[1]$.
The smooth-lisse proof above suffices for every conclusion below; no
additional assertion about arbitrary perverse restriction is needed.

\begin{theorem}[SG-09: the affine middle-degree bound]
\label{thm:affine}
Let $U/\mathbf F_q$ be smooth affine of pure dimension $d\ge1$, and
$\mathcal F$ lisse.  Suppose that, for a fixed $C<\infty$ and real $w$,
\[
 |\iota\operatorname{Tr}(F_{u,n},\mathcal F_{\bar u})|
       \le Cq^{nw/2}
 \quad\text{for all }n\ge1,\quad u\in U(\mathbf F_{q^n}).
\]
Every eigenvalue $\alpha$ on $H_c^d(U,\mathcal F)$ then satisfies
$|\iota\alpha|\le q^{(w+d)/2}$.
\end{theorem}
\begin{proof}
The empty case is immediate.  Choose a finite Noether normalization
$\pi:U\to\mathbf A^d$ \cite[Lemma 10.115.4, Tag 00OY]{StacksNN}.
Choose $B$ module generators of $\mathcal O(U)$ over the polynomial
coordinate ring.  Every geometric fiber has length at most $B$, hence
at most $B$ rational points over any finite extension.  Neither
\'etaleness nor constant fiber length is needed.
Choose affine coordinate generators $z_1,\ldots,z_s$ giving a closed
immersion $U\hookrightarrow\mathbf A^s$.

For each $k\ge2$, fix a nondegenerate bilinear form $b$ on
$\mathbf A^d$, and put
\[
 f(u_1,\ldots,u_k)=\sum_{i=1}^{k-1}b(\pi u_i,\pi u_{i+1}).
\]
Vary the per-vertex linear combinations $g_i=\sum_j a_{ij}z_j$ and $c$,
and consider
\begin{equation}\label{eq:pathsection}
 Y:\quad f+\sum_i g_i+c=0\quad\subset U^k.
\end{equation}
The map given by all the vertex coordinates is a closed immersion and
therefore quasi-finite.  Katz's Corollary 3.4.1, applied to this map,
the \emph{fixed} $f$, and $\mathcal F_k^\vee$, gives a dense open set of
coefficients with the restriction property of Proposition~\ref{prop:gysin}.
Generic smoothness of these hypersurfaces follows from hyperplane
Bertini for the graph closed immersion
$U^k\hookrightarrow\mathbf A^{sk+1}$, $u\mapsto(z(u),f(u))$.
The hyperplanes whose coefficient of the last coordinate is nonzero are
an open chart; rescale that coefficient to one.  Thus the smooth-good
open meets precisely the displayed family.  This argument does not
intersect a previously obtained open with an arbitrary possibly bad
fixed-coefficient slice.

Choose a point in the two good opens over some finite extension
$\mathbf F_{q^{e_k}}$.  It is permissible for $e_k$ and the coefficients
to depend on $k$.  Fix them, and write $Q=q^{e_kn}$ for any subsequent
extension degree $n\ge1$.  For $h\in\mathbf F_Q^*$ group the character
sum by $x_i=\pi u_i$, using the vertex weights
\[
 a_{i,h,n}(x)=\sum_{u\in\pi^{-1}(x)(\mathbf F_Q)}
       \iota\operatorname{Tr}(F_{u,e_kn},\mathcal F_{\bar u})
                         \psi(hg_i(u)).
\]
They have absolute value at most $BC Q^{w/2}$.  Lemma~\ref{lem:path}
applies although these weights depend on $i$ and $h$.  The zero mode
is exactly $M_{e_kn}^k$, since all vertex phases are then one.
Consequently
\begin{equation}\label{eq:parsevalbound}
 |\iota R_{k,e_kn}|\le(BC)^k
       q^{e_kn\{k(w+d)/2+d/2+1\}}.
\end{equation}
By Proposition~\ref{prop:gysin}, this remainder, up to its one sign,
is an effective power sum containing $\alpha^{e_kk}$.
Lemma~\ref{lem:power}, with $k$ fixed and all $n\ge1$, gives
\[
 |\iota\alpha|^{e_kk}
       \le q^{e_k\{k(w+d)/2+d/2+1\}}.
\]
Taking the $e_kk$th root and letting $k\to\infty$ proves the result.
No uniform bound on $e_k$ or the Betti numbers of $Y$ is required.
\end{proof}

\subsection{The weight consequences and their limits}

\begin{corollary}[SG-11: compact-support upper bounds]
\label{cor:upper}
Let $X/\mathbf F_q$ be separated of finite type, and $\mathcal F$ a
constructible $\overline{\mathbf Q}_\ell$-sheaf whose local Frobenius
eigenvalues are of $\iota$-absolute value at most $|\kappa(x)|^{w/2}$.
Then every eigenvalue on $H_c^i(X,\mathcal F)$ has $\iota$-absolute
value at most $q^{(w+i)/2}$.
\end{corollary}
\begin{proof}
First suppose $X$ is smooth affine of pure dimension $d$, with
$\mathcal F$ lisse.  In dimension zero the assertion is the local
hypothesis: induction of a residue-field Frobenius module replaces
each eigenvalue by its degree-many roots, without changing its weight.
For $d>0$, the groups in degrees $i<d$ vanish.  Theorem~\ref{thm:affine}
settles $i=d$, since the trace bound holds with the rank as $C$.
Choose another smooth generic hyperplane section $Y$; a finite base
extension is harmless.  For $i\ge d+1$ the Gysin map
\[
 H_c^{i-2}(Y,\mathcal F|_Y)(-1)\longrightarrow H_c^i(X,\mathcal F)
\]
is surjective at $i=d+1$ and isomorphic for $i>d+1$.
Induction on dimension bounds its source eigenvalues by
$q\,q^{(w+i-2)/2}=q^{(w+i)/2}$.
For a general $X$, remove nilpotents and choose a finite stratification
by smooth affine locally closed schemes on which $\mathcal F$ is lisse;
the finite field is perfect.  The localization exact sequence
\[
 H_c^i(U,\mathcal F|_U)\longrightarrow H_c^i(X,\mathcal F)
                   \longrightarrow H_c^i(Z,\mathcal F|_Z)
\]
for an open stratum union $U$ and its closed complement $Z$ shows,
by induction on the number of strata, that every middle eigenvalue
occurs in one of the adjacent groups.  Both have the desired bound.
The long exact sequence is used at degree $i$ on both sides, so there
is no extra shift loss.  A bound proved after a finite extension
descends by taking roots of the absolute values.
\end{proof}

This is the numerical weight inequality for compact support over a point,
as in \cite[Theorem 3.3.1]{WeilII}.  It is not a claim to have proved
algebraicity of arbitrary sheaf eigenvalues, all six-functor theorems,
or semisimplicity.  If the input bounds hold for every embedding under
consideration, the conclusion holds for each of them.

\begin{corollary}[SG-10: smooth proper purity and the requested Weil bounds]
\label{cor:purity}
Let $X/\mathbf F_q$ be smooth proper of pure dimension $d$, and let
$\mathcal F$ be pointwise $\iota$-pure of weight $w$ and lisse.
Then the eigenvalues on $H^i(X,\mathcal F)$ have absolute value
$q^{(w+i)/2}$.  In particular the finite-field absolute-value assertion
of Weil I holds for every smooth projective $X$, with constant
coefficients, without a Hodge standard conjecture.
\end{corollary}
\begin{proof}
Corollary~\ref{cor:upper} and properness give the upper bound.
The dual sheaf $\mathcal F^\vee$ is pointwise pure of weight $-w$.
Poincar\'e duality pairs
\[
 H^i(X,\mathcal F)\otimes
 H^{2d-i}(X,\mathcal F^\vee)\longrightarrow
                         \overline{\mathbf Q}_\ell(-d)
\]
perfectly.  Thus for every eigenvalue $\alpha$ in the first group,
$q^d/\alpha$ occurs in the second.  The latter's upper bound is
$q^{(-w+2d-i)/2}$, which gives the matching lower bound
$|\iota\alpha|\ge q^{(w+i)/2}$.
For constant coefficients the pointwise input is one, for every
embedding.  The rationality and integrality parts of Weil I are
separate established results, not deductions from the complex power-sum
lemma; the assertion reconstructed here is the requested absolute value.
\end{proof}

There is also a shorter projective-only conclusion from
Theorem~\ref{thm:affine}: choose a smooth hyperplane $H\subset X$,
use projective weak Lefschetz for $i<d$, and
$H_c^d(X\setminus H,\mathcal F)\to H^d(X,\mathcal F)\to
H^d(H,\mathcal F|_H)$ in the middle degree.  Induction gives the upper
bound, and duality gives the lower bound.  Both versions have been
checked; the compact-support proof records the stronger numerical scope
of the same mechanism.

\subsection{Why smoothness and spectral data alone are insufficient}

\begin{proposition}[HS-04 and HS-04B: the required genericity]
Smoothness of a hypersurface does not concentrate its compactly supported
cohomology in its middle degree.  Even with base middle concentration,
smoothness alone need not give the effective remainder of
Proposition~\ref{prop:gysin}.
\end{proposition}
\begin{proof}
For the first assertion take $U=\mathbf A^1$, $k=2$, constant weights,
and $Y:xy=1$.  Over $\mathbf F_{q^n}$, put $r=q^n$.
Then $M_n=r$, $N_n=r-1$, $R_n=-r$; the common top root in
$rN_n$ and $M_n^2$ cancels.  Indeed
$H_c^1(\mathbf G_m)=\overline{\mathbf Q}_\ell$ and
$H_c^2(\mathbf G_m)=\overline{\mathbf Q}_\ell(-1)$.
This example is compatible with Proposition~\ref{prop:gysin}, whose
base middle group is zero; it refutes the stronger concentration claim.

For the second assertion let $q$ be odd, let $\chi_n$ be the compatible
quadratic character of $\mathbf F_{q^n}$, extended by zero, and put
$t_n(x)=\chi_n(x^2-1)$.  This is the trace of the quadratic Kummer sheaf
on $U=\mathbf A^1\setminus\{1,-1\}$, extended by zero to the line.
The finite double cover $p:\mathbf G_m\to\mathbf A^1$,
$x=(u+u^{-1})/2$, has deck involution $u\mapsto u^{-1}$.
Its direct image splits as the constant sheaf and this extended Kummer
sheaf: at the two branch points the anti-invariant stalk is zero.
Subtracting the constant $H_c^2(\mathbf A^1)$ from the cohomology of
$\mathbf G_m$ shows that the Kummer sheaf has just
$H_c^1=\overline{\mathbf Q}_\ell$, with Frobenius eigenvalue one.
Thus $M_n=-1$ and the base is concentrated in its middle degree.

On the smooth curve $xy=1$ in $U^2$ one has
\[
 (x^2-1)(x^{-2}-1)=-\frac{(x^2-1)^2}{x^2}.
\]
Writing $\epsilon=\chi_1(-1)$ gives
\[
 N_n=\epsilon^n(q^n-3),\qquad
 R_n=(\epsilon q^2)^n-3(\epsilon q)^n-1.
\]
The expected sign $(-1)^{2-1}$ makes the coefficient of the root
$\epsilon q^2$ equal to $-1$.  Distinct finite power sums have unique
coefficients, by their generating-function poles, so $-R_n$ is not
effective.  Along this special curve the two ramification factors
cancel, leaving a nonzero top compact cohomology group.  The generic
restriction hypothesis fails.  Generic weak Lefschetz in the actual
Fourier-compatible family, rather than smoothness by itself, excludes
precisely this defect.
\end{proof}

\begin{proposition}[HS-05: weights and parity do not specify a polarization]
The Frobenius spectrum, its cohomological parity, and a Frobenius-compatible
nondegenerate symmetric form do not by themselves imply that form is
positive definite.
\end{proposition}
\begin{proof}
Place $V=\mathbf Q^2$ in degree $2r$ and let $F=q^r I$.
The forms with matrices $\operatorname{diag}(1,1)$ and
$\operatorname{diag}(1,-1)$ have the same compatibility
$b(Fx,Fy)=q^{2r}b(x,y)$ and identical spectral data, but signatures
$(2,0)$ and $(1,1)$.  This is a countermodel to the formal inference,
not a geometric counterexample to the Hodge standard conjecture.
Likewise a nontrivial Jordan block and its scalar semisimplification
have the same traces of all powers; the power-sum proof contains no
semisimplicity conclusion.
\end{proof}

\subsection{What was located in the supplied and retrieved Lean sources}
\label{sec:lean}
\paragraph{HS-06 (source-level conclusion).}
The short prose description is in Alp\"oge's \emph{Remark}, p.~3,
footnote 3, referring to timeline steps 3 and 5.  The supplied
\texttt{ALPOGE\_Main.lean} and \texttt{ALPOGE\_Statements.lean} are
wrappers for the stipulated global RH/Dirichlet-GRH statements.
For the method comparison we subsequently obtained the public release
archive specified in \texttt{evidence/LEAN\_MIRROR\_PROVENANCE.txt}; both
wrappers are byte-identical to their supplied versions.

The relevant retrieved file
\path{Development/EtaleCohomology/WeilConjectures/FermatCount.lean}
defines, in namespace \texttt{LateralTheory.EWeil}, the propositions
\texttt{Base}, \texttt{Goal}, \texttt{WeakAt}, \texttt{Tensor},
\texttt{FermatCount}, \texttt{HypTrace}, and others.
\texttt{Goal} has an explicit \texttt{Base} premise and concludes
$\|\iota\alpha\|^2=q^n$ for smooth projective cohomology;
\texttt{WeakAt} allows the interval $[n-3,n+3]$, and
\texttt{Tensor} sharpens such a dimension-independent loss.
The occurrence of a proposition definition or an implication between
propositions is not itself a closed theorem supplying all premises.
A targeted textual search of all 175,187 source files, and closer
declaration-level inspection of the 24 files in this directory and
the named divisibility/slope directories, are documented separately.
This is not a Lean kernel certificate or a proof of absence of a differently
named theorem.  The mathematical proof in this paper stands on its
explicit geometric inputs; no unlocated formal theorem is used to fill
a gap.  The arithmetic section gives exact declaration locators for the
steps 10 and 13 statements actually inspected.


\section{The sign, the evanescent part, and the Hodge-index route}

The sign-only construction of SG-08--SG-11 proves the required finite-field absolute values. Its decisive sign is the parity of an \emph{effective Frobenius remainder}: generic weak Lefschetz makes a Gysin map surjective, and the remaining virtual class is the class of an actual kernel together with a tensor power. This does not identify that sign with positive definiteness of an intersection form. Both arguments use hyperplane sections and evanescent cohomology, but the effective-remainder argument avoids the positive involution on correspondences required by the index route.

The supplied cote~16 scans~034, 036, 038 and~040 describe restriction, Gysin, Lefschetz operators, the evanescent summand, and its projector. They do not state positivity on the whole evanescent space. Scan~049 explicitly distinguishes algebraicity of an inverse from the extra index relation $I(X\times X)$ used for splittings and semisimplicity. Scan~050 announces a deduction of Weil from $C$ and $A$. The earlier definitions of those two labels are absent from the supplied excerpts, so their identification with modern standardized letters is not silently assumed. The conclusion here is precise: the sign-only proof realizes a cohomological hyperplane-induction programme, but it does not prove the additional cycle-theoretic hypotheses of the positive-involution argument announced on scans~049--050. The unrelated printed material on scans~035, 037 and~039 is an inventory fact, not by itself a transcription error.

\begin{proposition}[SG-01: parity is not intersection signature]
The sign $(-1)^j$ of a cohomology group in a trace formula, or its exponent $(-1)^{j+1}$ in the cohomological zeta function, does not determine the raw intersection sign on primitive algebraic classes in that degree. Neither sign gives a definite intersection form on the whole evanescent part of a hyperplane section.
\end{proposition}

\begin{proof}
Let $Q^{2r}$ be a split smooth even-dimensional quadric, as a smooth hyperplane section of a split $Q^{2r+1}$. Let $a,b$ be its two classes of maximal linear $r$-planes and put $\zeta=a-b$. The explicit calculation in SG-02 gives
\[
 h\zeta=0,\qquad \deg(\zeta^2)=2(-1)^r,
 \qquad E^{2r}\simeq\mathbb Q_\ell(-r).
\]
Here $E^{2r}$ is the kernel of the Gysin map
\[
 H^{2r}(Q^{2r}_{\overline{\mathbb F}_q},\mathbb Q_\ell)
 \longrightarrow H^{2r+2}(Q^{2r+1}_{\overline{\mathbb F}_q},\mathbb Q_\ell(1)).
\]
Its trace contribution over $\mathbb F_{q^n}$ is $+q^{rn}$, and its determinant occurs with exponent $-1$ in the zeta function. These two signs are independent of $r$, whereas $\zeta^2$ alternates with $r$. Thus $Q^2,Q^4,Q^6$ have trace signs $+,+,+$, determinant exponents $-,-,-$, and primitive squares $-2,+2,-2$. A uniform identification of parity with the raw index sign is false.

After normalizing the primitive form to $(-1)^r\deg(xy)$, all three Tate lines are positive. That numerical agreement with a positive multiplicity does not identify the constructions: one is a bilinear inequality on rational algebraic classes; the other is a sign in the Grothendieck group of Frobenius modules, defined without a pairing.

There is also no definite intersection form on the entire evanescent space in general. Embed $X=\mathbb P^2_{\mathbb C}$ by $\mathcal O(3)$, and take a smooth plane cubic $Y$, now a hyperplane section. Since $H^1(X)=H^3(X)=0$, its evanescent part is $E^1=H^1(Y,\mathbb Q)$, of dimension two. Its cup-product pairing is alternating by graded commutativity, and is nondegenerate by duality. Every vector has square zero; this cannot be a positive- or negative-definite real symmetric form. Its cohomological trace sign is nevertheless uniformly negative. This example concerns the whole evanescent space and does not contradict the Hodge standard conjecture, which concerns primitive \emph{algebraic} classes in even degree with a prescribed normalization.
\end{proof}

The exact cohomological locators are \cite[eq.~(1.5.4), p.~276]{WeilI} for determinant parity, \cite[\S4.1, p.~287]{WeilI} for the local vanishing-cycle squares $2,0,-2,0$ in dimensions congruent to $0,1,2,3$ modulo~$4$, and \cite[\S5.9, p.~293]{WeilI} for the alternating or symmetric form on $E/(E\cap E^\perp)$. The local Picard--Lefschetz signs are genuine intersection signs. They still do not assert positive definiteness on all of $E$.

\subsection{The split sixfold quadric and its Tate factor}

\begin{proposition}[SG-02: the split-quadric calculation]
Let $r\ge1$. On a split smooth $Q^{2r}$, the primitive middle cohomology and its corresponding rational Chow-motive summand are respectively $\mathbb Q_\ell(-r)$ and $\mathbb Q(-r)$. If $\zeta$ is the difference of the two maximal-plane classes, then $\deg(\zeta^2)=2(-1)^r$. Moreover,
\begin{align*}
 \#Q^{2r}(\mathbb F_{q^n})&=\sum_{j=0}^{2r}q^{jn}+q^{rn},\\
 Z(Q^{2r}/\mathbb F_q,T)&=
 \left((1-q^rT)\prod_{j=0}^{2r}(1-q^jT)\right)^{-1}.
\end{align*}
For the split sixfold model over $\mathbb Z$, the primitive global factor is $\zeta(s-3)$. Under the stipulated Alp\"oge Riemann hypothesis its nontrivial zeros have real part $7/2$. This is distinct from the local statement that its degree-six Frobenius eigenvalue is $q^3$.
\end{proposition}

\begin{proof}
Write the equation as $\sum_{i=0}^r x_i y_i=0$. The chart $x_0\ne0$ is $\mathbb A^{2r}$. Its complement is the projective cone on $Q^{2r-2}$: it has one vertex, and its complement of the vertex is an $\mathbb A^1$-bundle over $Q^{2r-2}$. Starting with $Q^0$, which consists of two points, this gives an affine paving with one cell in every dimension $0,1,\ldots,2r$ and an extra cell in dimension $r$. All these cells are defined over the split base. Counting the cells proves the point-count formula, and exponentiating $\sum_{n\ge1}\#Q(\mathbb F_{q^n})T^n/n$ proves the zeta formula. The cell filtration and the localization sequence give only even cohomology, with one Tate class in every even degree and an additional middle class. The two maximal-plane classes below form an explicit Frobenius eigenbasis in middle degree, so no semisimplicity or weight theorem is being used to split a possible extension.

For a maximal plane $\mathbb P^r\subset Q^{2r}$ the normal exact sequence is
\[
 0\longrightarrow N_{\mathbb P^r/Q^{2r}}
 \longrightarrow\mathcal O(1)^{r+1}
 \longrightarrow\mathcal O(2)\longrightarrow0.
\]
Consequently its self-intersection is the coefficient of $t^r$ in
$(1+t)^{r+1}/(1+2t)$, namely $(1+(-1)^r)/2$. A special linear
$\mathbb P^{r+1}$ section is the union of a plane from each ruling, so
$h^r=a+b$. Since $\deg(h^r a)=\deg(h^r b)=1$, it follows that
\[
 a^2=b^2=\frac{1+(-1)^r}{2},\qquad
 ab=\frac{1-(-1)^r}{2},\qquad
 (a-b)^2=2(-1)^r.
\]
The intersection matrix is nonsingular, proving the independence of the two classes. In codimension $r+1$ the cell calculation gives a single generator, and $ha=hb$ by pairing with $h^{r-1}$; hence $h\zeta=0$. The same equality holds in the Chow ring: the cell closures generate its groups, and their cycle classes, of the same rank, are independent. The correspondence
\[
 p_{\mathrm{prim}}=\frac{\zeta\times\zeta}{\deg(\zeta^2)}
\]
acts identically on $\zeta$ and kills the hyperplane line. The cycle $\zeta$ and its transposed cycle divided by $\deg(\zeta^2)$ give inverse maps between this summand and $\mathbb Q(-r)$. This proves the assertion about the rational Chow motive. For $Q^{2r}\subset Q^{2r+1}$, the ambient middle class has nonzero Gysin image and $\zeta$ has zero Gysin image. Since the kernel has dimension one, it is exactly this primitive Tate line.

For $r=3$, the model $\sum_{i=0}^3x_i y_i=0$ in $\mathbb P^7_{\mathbb Z}$ is smooth over $\mathbb Z$: its partial derivatives are the coordinates $x_i,y_i$, which cannot vanish simultaneously at a projective point. Every finite-field fibre is split. Its primitive Euler factor is $(1-p^{3-s})^{-1}$, so their product is $\zeta(s-3)$ in the half-plane of absolute convergence, and then by meromorphic continuation. The change of variable $u=s-3$ carries the stipulated Riemann hypothesis to $\operatorname{Re}(s)=7/2$.
\end{proof}

Thus the requested absolute values for the ambient quadric have a direct Tate-motive computation; SG-08--SG-11 also give them through the general sign-only theorem. The elementary intersection inequality is independently visible as $(-1)^r\zeta^2=2>0$. Neither computation should be confused with a deduction of the Hodge standard conjecture from trace parity.

\subsection{The extra datum in the positive-involution proof}

\begin{lemma}[SG-03: the exact positivity mechanism]
Let $A$ be a finite-dimensional real algebra with identity, $\dagger$ an involutive anti-automorphism, and $\tau:A\to\mathbb R$ a cyclic linear functional satisfying $\tau(a^\dagger)=\tau(a)$. Suppose
\[
 B(a,b)=\tau(ab^\dagger)
\]
is positive definite. If $F^\dagger F=c\,1$ with $c>0$, every eigenvalue of $F$ in a faithful finite-dimensional representation has modulus $\sqrt c$.
\end{lemma}

\begin{proof}
The equation makes $F$ invertible, hence $FF^\dagger=c$ also. Cyclicity gives
\[
 B(Fa,Fb)=\tau(Fab^\dagger F^\dagger)
        =\tau(F^\dagger F ab^\dagger)=cB(a,b).
\]
Left multiplication by $c^{-1/2}F$ is therefore an isometry of a positive-definite real vector space. Its complex eigenvalues have modulus one. The minimum polynomial of left multiplication by $F$ equals the minimum polynomial of $F$ as an element of $A$: indeed $p(L_F)=L_{p(F)}$, and applying to $1$ detects zero. A faithful representation has that same minimum polynomial, proving the assertion.
\end{proof}

This short lemma states the precise form needed here of the classical positive-Rosati argument; compare \cite[Th.~13.2 and Lemmas~14.2--14.3]{MilneAV}. In the correspondence-algebra route, construction of the positive form and its algebraic adjoint is additional content. Merely knowing an effective trace remainder does not produce it. For a sharp formal comparison, the indefinite matrix
\[
 J=\begin{pmatrix}0&1\\1&0\end{pmatrix},\qquad
 F=\operatorname{diag}(\sqrt c\,a,\sqrt c/a),\quad a>1,
\]
satisfies $F^tJF=cJ$ but has eigenvalues of different absolute values. Conversely $F=\sqrt c\,I$ has pure absolute values while preserving this indefinite form up to the same multiplier. These are countermodels between linear algebra data, not geometric counterexamples to the Hodge standard conjecture. Cote~16 scan~049 is precisely where the extra index input is acknowledged.

\subsection{CM elliptic powers and arbitrary abelian varieties}

\begin{proposition}[SG-04: the precise hexagonal-curve case]
Let $E/\mathbb C$ be the elliptic curve $\mathbb C/\mathbb Z[\omega]$. Every rational algebraic cohomology class on $E^n$ is generated by divisor classes. The same assertion holds for the geometric rational algebraic $\ell$-adic classes on powers of every good finite-field reduction of $E$. The Hodge standard conjecture on these class spaces is consequently an unconditional instance of \cite[Remarks~3.7--3.8]{MilnePolarizations}. This concerns cohomology classes, and makes no assertion that every Chow class or every class modulo algebraic equivalence is a polynomial in divisors.
\end{proposition}

\begin{proof}
Over $\mathbb C$, put $K=\mathbb Q(\omega)$. Then
$\operatorname{End}^0(E^n)=M_n(K)$, and Rosati for the product polarization is conjugate transpose. Its fixed subspace has rational dimension
$n+2\binom n2=n^2$. The N\'eron--Severi/Rosati identification
\cite[Prop.~13.1]{MilneAV} therefore gives $\dim_{\mathbb Q}\operatorname{NS}(E^n)_{\mathbb Q}=n^2$. The first Chern class is injective on this space. Since $H^{1,1}(E^n)$ has complex dimension $n^2$, divisor classes span it over $\mathbb C$.
Write $H^1(E^n,\mathbb C)=U\oplus W$ in Hodge types $(1,0),(0,1)$, each of dimension $n$. The space
\[
 H^{r,r}(E^n)=\bigwedge^rU\otimes\bigwedge^rW
\]
is spanned by products $(u_1\wedge w_1)\cdots(u_r\wedge w_r)$ of degree-two classes. Thus the complexification of the rational algebra generated by divisors fills $H^{r,r}$. Every rational algebraic class belongs to that rational algebra. Here we use the exterior-algebra description of abelian cohomology, \cite[Th.~11.1(b), and Th.~2.3 for Betti cohomology]{MilneAV}.

For completeness the finite-field assertion requires both geometric endomorphisms and rational descent of coefficients. The ordinary/supersingular classification and the criterion $p\nmid\operatorname{Tr}(\pi_E)$ in the ordinary case are \cite[Th.~13.3, Th.~13.18, Cor.~13.20]{SutherlandEC}. For an ordinary elliptic curve, the two roots $\alpha,\beta$ of the Frobenius polynomial satisfy $\alpha\beta=q$ and $\alpha/\beta$ is not a root of unity. Indeed they are algebraic integers, their sum is a $p$-adic unit and their product has positive valuation, so one has valuation zero and the other has valuation $v_p(q)$. A geometric algebraic class is fixed after a finite extension and a Tate twist. In degree $2r$ on $E^n$, the invariant monomials therefore have exactly $r$ vectors from each of the two eigenspaces. They are products of paired degree-two classes. The geometric endomorphism algebra is imaginary quadratic, so the Rosati calculation above gives N\'eron--Severi rank $n^2$, exactly the dimension of the paired degree-two space. The divisor cycle map fills that space.

For a supersingular elliptic curve the geometric endomorphism algebra is quaternionic and Rosati is its canonical conjugation. The self-adjoint matrices in $M_n(\operatorname{End}^0(E))$ have dimension
\[
 n+4\binom n2=n(2n-1)=\binom{2n}{2}.
\]
Thus divisor classes span all of $H^2$. The even exterior algebra on $H^1$ is generated by $H^2$, so divisor products span every even cohomology group over $\mathbb Q_\ell$.

To upgrade these spanning assertions to \emph{rational} algebraic classes, use Poincar\'e duality. It is perfect on the balanced monomial spaces in the ordinary case, and on all even cohomology in the supersingular case. Choose divisor-product bases $b_i,d_j$ of complementary spaces over $\mathbb Q_\ell$. Their intersection matrix $M=(\deg(b_i d_j))$ has rational entries and is invertible. If an algebraic cycle $z$ has
$\operatorname{cl}(z)=\sum_i c_i\operatorname{cl}(b_i)$, then its pairings with the $d_j$ give an invertible rational linear system with rational right-hand side. Hence all $c_i$ are rational. This proves the precise generation assertion.

Finally, \cite[Remark~3.7]{MilnePolarizations} proves that
$(-1)^r\deg(L^{n-2r}xy)$ is positive definite on primitive Lefschetz classes, for $0\le r\le n/2$, over every algebraically closed field. The class spaces just identified are all algebraic cohomology classes in the stated examples, giving the claimed instance of Hodge standard positivity.
\end{proof}

\begin{remark}
The complex argument uses only that $\operatorname{End}^0(E)$ is imaginary quadratic, so it also proves the same assertion for powers of any CM \emph{elliptic} curve. The finite-field argument applies to powers of every elliptic curve over a finite field. An arbitrary higher-genus curve with CM Jacobian, or an arbitrary CM abelian variety, is not covered by that argument. For general abelian varieties, \cite[Th.~3.3]{MilnePolarizations} is conditional on the Hodge conjecture for all complex CM abelian varieties; its unconditional Remark~3.7 concerns Lefschetz classes. No unspecified Masullo statement is substituted for that hypothesis.
\end{remark}

\begin{corollary}[SG-05: all abelian absolute values]
For every abelian variety $A/\mathbb F_q$, every eigenvalue on $H^i(A_{\overline{\mathbb F}_q},\mathbb Q_\ell)$ has absolute value $q^{i/2}$. This is available independently of the full Hodge standard conjecture in all codimensions.
\end{corollary}

\begin{proof}
The classical inputs are Rosati positivity, $\pi^\dagger\pi=q$, and the resulting degree-one bound, precisely \cite[Th.~13.2, Lemmas~14.2--14.3 and Th.~14.1]{MilneAV}. The exterior-algebra isomorphism \cite[Th.~11.1(b)]{MilneAV} makes every degree-$i$ eigenvalue a product of $i$ degree-one eigenvalues. Their moduli multiply to $q^{i/2}$. This is a short deduction from the cited theorem, not a new proof of Rosati positivity. The main sign-only theorem supplies a second proof of these absolute values.
\end{proof}

Hence the positivity mechanism already works in degree one for every abelian variety, and for all primitive algebraic classes of the elliptic powers just treated. None of these facts identifies the Gysin-remainder sign with an index form, or proves a new general positivity theorem for exotic algebraic classes.

\subsection{The letter's absolute values and all SA06 C.1 rows}

\begin{proposition}[SG-06: fibre purity and the invisible cycle kernel]
Every smooth projective finite-field fibre of a smooth projective spread of a complex variety in the supplied SA06 C.1 rows has the absolute values $q^{i/2}$ in degree $i$, by SG-08--SG-11 (equivalently, the known \cite[Th.~1.6]{WeilI}). These absolute values do not recover the kernel of the cycle class map or upgrade statements formulated modulo algebraic equivalence.
\end{proposition}

\begin{proof}
A projective embedding and finitely many equations descend to a finitely generated integral subring of $\mathbb C$. After shrinking its spectrum to the smooth locus one obtains a smooth projective model $\mathcal X\to S$. At each closed point of $S$ the residue field is finite, and the main sign-only theorem applies to the corresponding smooth projective fibre. This is meaningful even when the original very general complex example requires a transcendental field of definition; there is no canonical single finite-field Frobenius on that original complex fibre.

If $z$ has zero cohomology class and $w$ is a complementary algebraic cycle, compatibility of cycle class with products and trace gives
\[
 \deg(z\cdot w)=\operatorname{Tr}\bigl(\operatorname{cl}(z)\cup\operatorname{cl}(w)\bigr)=0.
\]
Every cohomological correspondence operator also factors through its cycle class. Consequently its traces, characteristic polynomials and local $L$-factors are unchanged on replacing a cycle by one with the same cohomology class. Spectral information therefore does not distinguish a nonzero Griffiths class from zero, or establish that it vanishes modulo algebraic equivalence. This is a factorization statement, not a claim of a finite-field counterexample to conjectured rational vanishing of Griffiths groups.
\end{proof}

The supplied intake provides full-check descriptions for C.1a--e and only abbreviated statement descriptions for C.1f--q. The table records every row without claiming to have independently checked the absent original proofs. In each row, ``purity'' may now be obtained by the reconstructed sign-only argument, not merely by citing Deligne. The remaining distinctions concern what purity does not assert about cycles.

{\small
\begin{longtable}{@{}p{0.08\textwidth}p{0.34\textwidth}p{0.52\textwidth}@{}}
\toprule
Row & Supplied variety or assertion & Consequence for this question \\
\midrule
\endfirsthead
\toprule
Row & Supplied variety or assertion & Consequence for this question \\
\midrule
\endhead
C.1a & Very general quintic threefold $V$; infinite Griffiths summand in $A^2_{\mathrm{alg}}$. & Every smooth finite-field fibre is pure. Its Griffiths kernel is invisible to traces. Its varying $H^3$ is not identified with the Tate line $\mathbb Q(-3)$; no number-field $H^3$-GRH follows. \\
C.1b & Failure of the letter's $A,A',A''$ modulo algebraic equivalence. & Fibre purity does not remove a cycle-class-kernel obstruction in that formulation. \\
C.1c & $L^{n-2}:\operatorname{NS}_{\mathbb Q}\to A^{n-1}_{\mathrm{alg}}$; examples $V\times\mathbb P^{n-3}$. & K\"unneth gives $H^i(V_s\times\mathbb P^m)=\bigoplus_j H^{i-2j}(V_s)(-j)$, with the required weights. The stated Griffiths cokernel remains separate. \\
C.1d & Endpoint cases, dimensions at most two, and surface Hodge index. & Curve purity also has the classical Hodge-index/Rosati proof. Parity is still not its index form. \\
C.1e & $V\times\mathbb P^1$; vacuous primitive letter-$B$ in the stated degree, with a nonzero numerical kernel. & Product purity does not imply algebraic equivalence equals numerical equivalence. The row's vacuity qualification is retained; this is not a counterexample to primitive Hodge standard positivity. \\
C.1f & Universal ranges of letter-$A$ over $\mathbb C$. & Purity in every cohomological degree does not enlarge valid ranges of a statement modulo algebraic equivalence. \\
C.1g & Universal quintic over $\mathbb Q(t_1,\ldots,t_{113})$ with a homologically trivial cycle. & Use a finite-type $\mathbb Z$-model and its smooth finite-field closed fibres. A generic transcendental base has no single arithmetic Frobenius. \\
C.1h & Geometric generic hyperplane and its Gysin image modulo algebraic equivalence. & Smooth closed fibres are pure; the asserted Chow-level Gysin image needs its separate proof. \\
C.1i & Universal $A'/A''$ range and split-quadric examples. & SG-02 gives the exact Tate eigenvalues for the quadrics. Their purity cannot repair an incorrectly quantified restriction or Gysin assertion. \\
C.1j & Middle letter-$B$ and injectivity on the Griffiths kernel. & Positivity on cohomological realizations and injectivity on an algebraic-equivalence kernel are different requirements. \\
C.1k & Cycle-class injectivity, finite rank or generation, and equivalence relations. & Frobenius purity neither supplies injectivity nor controls the size of its kernel. \\
C.1l & Radical correspondences constructed from Griffiths classes. & A zero cohomology class acts as zero on cohomology; its spectral data vanish without proving it zero modulo algebraic equivalence. \\
C.1m & Semisimplicity modulo algebraic equivalence in universal dimension ranges. & Purity alone does not give general Frobenius semisimplicity or semisimplicity of this correspondence algebra. Cote~16:049 expressly invokes $I(X\times X)$. \\
C.1n & An endpoint image criterion, standard $D$, and a kernel condition. & Weight purity supplies neither that cycle-equivalence assertion nor the extra kernel vanishing. \\
C.1o & Projective-space stabilization and a middle-positivity radical. & Tate shifts preserve the required weights; stabilization does not make a zero cycle class visible to Frobenius. \\
C.1p & Abelian projective stabilizations and Beauville/Griffiths pieces. & SG-05 and projective Tate shifts give purity. They do not determine the Beauville pieces of a Griffiths kernel. \\
C.1q & Blow-up exceptional normal-bundle polynomial and a radical. & Smooth projective finite-field blow-ups are pure; any exceptional-cycle or radical assertion remains a Chow-level assertion. The intake's truncated description does not authorize a stronger reconstruction. \\
\bottomrule
\end{longtable}
}

The answer to the letter's absolute-value question is therefore affirmative for every applicable finite-field fibre, using the sign-only proof. The stronger claim that this restores the letter's formulations modulo algebraic equivalence is not supported and is contradicted, in their stated complex scope, by the prior-run counterexamples being kept separate. The stipulated Dirichlet GRH transfers to the ambient Tate factor $\zeta(s-3)$ by SG-02. It does not establish a number-field GRH for an unshown non-Artin--Tate factor such as a general quintic's $H^3$.

\subsection{The ambient quadric versus the E.3u hypersurface}

\begin{proposition}[SG-07: restriction kills the ambient primitive line]
Let $j:Y\hookrightarrow Q^6$ be a smooth hypersurface of positive degree $d$, including the degree-$14$ section in the supplied SA06 E.3u description. The primitive class $\zeta\in H^6(Q^6,\mathbb Q_\ell)$ satisfies $j^*\zeta=0$.
\end{proposition}

\begin{proof}
The projection formula gives
\[
 j_*j^*\zeta=d\,h\zeta=0.
\]
The complement $U=Q^6\setminus Y$ is affine, since $Y$ is an effective ample divisor. Artin vanishing and Poincar\'e duality give $H_c^a(U,\mathbb Q_\ell)=0$ for $a<6$. The localization sequence therefore identifies $H^4(Q^6)\to H^4(Y)$. Dualizing gives an isomorphism
\[
 j_*:H^6(Y,\mathbb Q_\ell)
 \longrightarrow H^8(Q^6,\mathbb Q_\ell)(1).
\]
It follows that $j^*\zeta=0$. The twist is essential: both sides have Frobenius eigenvalue $q^3$ on the indicated Tate summand. The vanishing input is SGA~4, XIV, Th.~3.1; this is the ordinary ample-section proof of weak Lefschetz and requires no genericity assumption on the given smooth $Y$.
\end{proof}

This explains why the difference of the two plane traces can become a homologically trivial cycle on $Y$. Its morphic or Lawson nontriviality is additional structure. The single ambient Tate summand cannot by itself supply the variable $H^5(Y)$, nor see the Griffiths or Lawson kernel on $Y\times E$. The general sign-only theorem nevertheless proves the Frobenius absolute values of every smooth projective finite-field fibre of those varieties. The degree-$14$ nonvanishing assertion attributed to Walker and Nagel in SA06 E.3u is not independently verified by this argument, and is not needed for it.


\section{Cote 44, arithmetic slopes, and integral Lefschetz operators}
\label{sec:arithmetic-sl2}

The conclusions in this section have three distinct statuses.  The literal
height comparison in cote 44, Theorem 5, has a geometric counterexample under
all its explicitly transcribed hypotheses.  Several elementary integral
Lefschetz statements admit complete proofs and exact matrix checks.  A
comparison between an actual theta-lift chain and those Lefschetz lattices
remains an additional mathematical assertion.  The accepted RH/GRH input is
not questioned or audited here.

\subsection{AR-01: a counterexample to the displayed higher-cycle height formula}
\label{ar:counterexample}

\paragraph{Statement.}
Over an algebraically closed field $k$ of characteristic zero there is a
smooth projective fourfold $X$ with a Lefschetz pencil of hyperplane sections,
all of whose fibres are geometrically integral, and an integral class
$Z\in\mathrm{CH}^2(X)$ such that
\[
 HZ=0,\qquad Z|_Y=0\text{ in }\mathrm{CH}^2(Y),\qquad \deg Z^2=3,
\]
where $Y$ is the generic member and $H$ is the hyperplane class.  Consequently
the displayed assertion of cote 44, scan 004, Theorem 5,
\[
 \langle Z|_Y,Z|_Y\rangle=(-1)^m\deg Z^2
\]
is false for $m=2$ under its explicitly written assumptions.  This also
refutes H9(a)'s assertion that integral fibres alone remove the need for
vertical corrections in higher codimension.

\begin{proof}
Let $X=\mathbb P^2_k\times\mathbb P^2_k$, with hyperplane classes $a,b$,
and embed it by the Segre polarization $H=a+b$.  Its Chow ring is
\[
 \mathrm{CH}^*(X)=\mathbb Z[a,b]/(a^3,b^3),\qquad \deg(a^2b^2)=1.
\]
For $Z=a^2-ab+b^2$ one computes
\[
 HZ=a^3+b^3=0,\qquad
 Z^2=a^4-2a^3b+3a^2b^2-2ab^3+b^4=3a^2b^2.
\]
Consider the pencil
\[
 \lambda(x_0y_0+x_1y_1+x_2y_2)
       +\mu(x_1y_1+2x_2y_2)=0.
\]
Its coefficient matrix is
$\operatorname{diag}(\lambda,\lambda+\mu,\lambda+2\mu)$.  It has rank
three except at $[0:1],[-1:1],[-2:1]$, where it has rank two.  A bilinear
polynomial of rank at least two is geometrically irreducible: a nontrivial
factorization of its bidegree $(1,1)$ must have factors of bidegrees $(1,0)$
and $(0,1)$, which would force rank one.  Every member is therefore
geometrically integral.

A rank-three member is smooth.  A rank-two diagonal member has exactly one
singular point, $(e_i,e_i)$, corresponding to the vanishing diagonal entry.
On the chart $x_i=y_i=1$, its equation is the sum of two nonzero multiples
of $x_jy_j$, $j\ne i$.  This is a nondegenerate quadratic form in four
variables, hence an ordinary double point.  None of the three singular
points lies on the base locus: at it $\sum x_jy_j=1$.

The two defining divisors have no common component, so the base is a
codimension-two complete intersection.  If its gradients were dependent
at a base point, that point would be singular on a member of the pencil,
contrary to the preceding description.  The base is smooth.  Its blow-up
therefore gives a Lefschetz pencil.  Each fibre is isomorphic to the
corresponding member, because the base is a Cartier divisor in that member.

Put $K=k(\mathbb P^1)$.  The generic coefficient matrix is invertible, so
a $K$-linear change of $y$ coordinates identifies $Y$ with
$\sum x_i y_i=0$.  Projection to $\mathbb P^2_x$ makes $Y$ the bundle of
lines in
\[
 E=\ker\bigl(\mathcal O^3\longrightarrow\mathcal O(1)\bigr),\qquad
 c(E)=1-a+a^2.
\]
The class $b$ is the first Chern class of the dual tautological line.
The projective bundle relation is
\[
 b^2+c_1(E)b+c_2(E)=b^2-ab+a^2=0.
\]
For clarity, this sign also follows directly from the tautological exact
sequence: its two line factors have first Chern classes $-b$ and $-a+b$,
so $a^2=c_2(E)=(-b)(-a+b)$.  Thus $Z|_Y=0$ in the Chow group.  Every
bilinear height pairing on the restricted classes sends $(0,0)$ to zero,
whereas the displayed formula would give $0=3$.
\end{proof}

The Chow calculation uses the usual projective bundle formula
\cite[Theorem 3.3(b)]{Fulton}; the particular relation is derived above.
This is a counterexample to the raw identity, not to the Hodge sign:
$Z$ is primitive and its square has the correct positive sign.  The generic
fibre is cellular, so its algebraic intermediate Jacobian in this degree
is zero.  In particular, an injective isometric realization of all primitive
middle cycles through that Jacobian is impossible.

The source's margin ends with an ellipsis after the statement that the
fibres are geometrically integral.  No claim is made that an untranscribed
admissibility condition could not have been intended.  In codimension
$m\ge2$, an integral fibre can still support codimension-$(m-1)$ cycles,
which become vertical codimension-$m$ cycles of the total space.  Fibre
integrality removes additional vertical divisors on a surface fibration;
it does not remove these higher-cycle correction groups.

\subsection{AR-02: the sign is absent from slope and norm data}
\label{ar:slope-sign}

\paragraph{Statement.}
A Hermitian lattice, an injective map to its Hermitian dual, all its local
operator norms, and Bost's slope inequality do not determine the sign of
the associated symmetric bilinear form.

\begin{proof}
Take $E=\mathbb Z$ with its standard metric and $E^\vee$ with the dual
metric.  The maps $u_+(1)=1$ and $u_-(1)=-1$ are integral isomorphisms of
norm one at every finite and infinite place.  Their degrees, maximal slopes
and map heights coincide and are zero; both satisfy the slope inequality
with equality.  Their symmetric forms are $xy$ and $-xy$, with opposite
signatures.  Multiplication of a map by $-1$ preserves all these norm data.
\end{proof}

Bost's precise statement for Hermitian bundles $E,F$ over $\mathcal O_K$
and an injective $K$-linear map $\varphi$ is
\[
 \widehat{\deg}_{n}(E)
 \leq\operatorname{rk}(E)
       \bigl(\widehat\mu_{\max}(F)+h(E,F,\varphi)\bigr).
\]
The locator is \cite[\S4.1.4, Proposition 4.5, equation (4.11)]{Bost};
the filtered version is Proposition 4.6, equation (4.14).  Their complete
printed proofs were read: the first uses the saturated image, determinant
norm bounds and maximal slope; the second applies it to filtered quotients
and sums degrees.  Basic arithmetic-degree properties are cited background.

\paragraph{The retrieved step 13 declaration.}
The actual files in \path{Formalism/LatticeSlopes} sharpen this scope.
In \texttt{Statement.lean}, \texttt{Formalism.F.SlopeInequality} concerns
discrete $\mathbb Z$-lattices in finite-dimensional positive Euclidean
spaces.  If $L\ne0$ and $\varphi:L\to M$ is an injective integral map
with $\|\varphi x\|\le C\|x\|$, it asserts
\[
 \mu(L)\le\mu_{\max}(M)+\log C,\qquad
 \deg(L)=-\log\operatorname{covol}(L),\quad
 \mu(L)=\deg(L)/\operatorname{rk}(L).
\]
The hypotheses imply $C>0$.  The complete proof body
\texttt{slopeInequality\_of} in
\path{LatticeSlopesProofs/SlopeInequality.lean} was read.  It proves
\[
 \operatorname{covol}(\varphi L)
 \le C^{\operatorname{rk}(L)}\operatorname{covol}(L)
\]
using a Gram determinant, and applies the maximum-slope bound to the image
lattice.  A positive inner product is an input, not a conclusion about an
algebraic-cycle intersection form.  This was source inspection, not a Lean
kernel run.

The separate file
\path{ReciprocityTheory/ReductiveGroups/ArithmeticGeometry/}
\path{ArakelovGeometry/MoriwakiSt.lean} defines
\texttt{MoriwakiHSStatement} and \texttt{MoriwakiBigStatement}.  These
are arithmetic Hilbert--Samuel asymptotics and a volume bound under
semipositivity and nefness hypotheses; the asymptotic statement also assumes
nonnegative heights.  Neither statement asserts a primitive intersection
signature.  Their full upstream proofs were not inspected.

\subsection{AR-03: metric dependence of an arithmetic self-intersection}
\label{ar:metric}

Let $X=\mathbb P^1_{\mathbb Z}$ and let $\overline L_0$ be
$\mathcal O(1)$ with a smooth conjugation-invariant metric.  Scale the
squared norm by $e^{-t}$ and call the result $\overline L_t$.  With
$x=\widehat c_1(\overline L_0)$ and the convention
$\widehat{\deg}a(\eta)=\frac12\int\eta$, one has
\[
 \widehat{\deg}\,\widehat c_1(\overline L_t)^2
       =\widehat{\deg}(x^2)+t.
\]
\begin{proof}
The representative $-\log\|s\|^2$ for the Green function gains the
constant $t$, so $\widehat c_1(\overline L_t)=x+a(t)$.  The arithmetic
product gives $x\,a(t)=a(t\,c_1(\overline L_0))$ and $a(t)^2=0$,
since $dd^ct=0$.  Expanding the square and using
$\int_{\mathbb P^1(\mathbb C)}c_1(\overline L_0)=1$ proves the formula.
\end{proof}
The product conventions are those of
\cite[\S4.2.11 and \S4.3, especially Theorem 4.3.2]{GilletSoule}.
Thus an arbitrary sign of this arithmetic self-intersection is compatible
with identical finite fibres and Frobenius data.  This does not contradict
a properly primitive arithmetic Hodge-index theorem with its ampleness and
metric hypotheses.

\subsection{AR-04: known cases and the remaining arithmetic comparison}
\label{ar:known}

Cote 44, Theorem 4, concerns the canonical self-duality of $P^m(X)$ for
$\dim X=2m-1$ over an algebraically closed field; the field hypothesis is
explicit in scan 008.  It is not by itself an Arakelov statement over
$\mathcal O_K$.  Theorem 5 concerns function-field heights in a pencil.
The primitive-decomposition margin in scan 008 must remain part of any
reconstruction of the sign convention.

The following are cited background, rather than new consequences of the
accepted RH/GRH theorem.
\begin{itemize}
\item For $m=1$, a curve Jacobian has its classical principal polarization
\cite[Theorem 6.6 and Summary 6.11]{MilneJacobian}.
For a cellular generic fibre in AR-01 the higher algebraic Jacobian is zero,
so its polarization assertion is vacuous.
\item Moriwaki's Theorem A(2) proves, for a regular projective flat
arithmetic variety of relative dimension $d\ge1$ with arithmetically ample
$\overline H$, that nonzero $x\in\widehat{\mathrm{CH}}^1(X)_{\mathbb R}$
satisfying $\widehat c_1(\overline H)^d x=0$ obeys
\[
 \widehat{\deg}\bigl(x^2\widehat c_1(\overline H)^{d-1}\bigr)<0.
\]
Theorem B gives the weak inequality for generic-degree-zero divisors and
characterizes equality by a positive multiple pulled back from
$\operatorname{Spec}\mathcal O_K$ \cite[Theorems A and B]{Moriwaki}.
The exact hypotheses and statements were checked; the complete proof was
not re-proved.
\item K\"unnemann proves the arithmetic standard conjectures for projective
space with the specified canonical Fubini--Study metric
\cite[Theorem 4.1]{Kunnemann}.  His discussion of abelian schemes keeps
the positivity of heights of higher homologically trivial cycles as an
additional conjectural input.  It must not be reported as a theorem for all
higher cycles on abelian varieties.
\item In surface pencils, divisor-height identities include vertical
component corrections.  For the elliptic case, the exact primary locators
are \cite[Lemma 8.1, Theorems 8.4 and 8.6, equations (8.8), (8.11)--(8.14)]{Shioda}.
The correction coefficients come from the negative inverse component
intersection matrix.  These statements were checked at the primary-source
text level, not re-proved in full.
\end{itemize}

No new higher polarization or arithmetic Hodge-index case follows from the
two bare steps 5 and 13.  This is a statement about the missing derivation,
not a claim that every possible stronger comparison is false.  A corrected
higher version needs admissible cycle extensions, vertical correction
groups, and an identity with the particular canonical self-duality and its
primitive signs.  AR-01 proves that the uncorrected universal identity
cannot supply that identity.

\subsection{AR-05: a precise positive consequence of divisibility and slopes}
\label{ar:divisibility-slope}

Let $L,M$ be finite-rank discrete lattices in positive Euclidean spaces,
with $L\ne0$.  If $\varphi:L\to M$ is injective, $\|\varphi\|\le C$
with $C>0$, and
$\varphi L\subset dM$ for a positive integer $d$, then
\[
 \log d\le\mu_{\max}(M)-\mu(L)+\log C.
\]
\begin{proof}
The map $\varphi/d$ is integral and injective, with norm at most $C/d$.
Apply the inspected slope inequality and rearrange.
\end{proof}
If $L$ and $M$ have the same rank $r$, one has the sharper estimate
\[
 \log[M:\varphi L]\le r\log C+\deg(M)-\deg(L).
\]
\begin{proof}
The extension to real spans has absolute determinant equal to the product
of its singular values, hence at most $C^r$.  Therefore
$\operatorname{covol}(\varphi L)\le C^r\operatorname{covol}(L)$.
Use $\operatorname{covol}(\varphi L)=[M:\varphi L]\operatorname{covol}(M)$
and take logarithms.  If $\varphi L\subset dM$, then
$[M:\varphi L]\ge d^r$, giving the corresponding bound with $\mu(M)$.
\end{proof}
This explains how an arithmetic divisibility lower bound can compete with
an archimedean norm upper bound.  It still does not determine a signature:
the maps $\pm d\,\mathrm{id}$ on $\mathbb Z^r$ have the same exact
index $d^r$, norm $d$ and slopes, with equality in these estimates, but
opposite associated symmetric forms.

\subsection{SL-01: integral blocks for divided-power Lefschetz maps}
\label{sl:blocks}

Let $M=\bigoplus_{a=1}^n(\mathbb Ze_a\oplus\mathbb Zf_a)$,
$p_a=e_a\wedge f_a$, $\xi=\sum_a p_a$, and $L(u)=\xi\wedge u$.
The divided power $L^{(m)}=L^m/m!$ is integral because
\[
 \xi^m/m!=\sum_{|R|=m}p_R.
\]
Fix $s$ pairs on which a basis wedge uses exactly one generator, and choose
which generator in each.  Set $v=n-s$.  The degree-$j=s+2t$ block is
indexed by $t$-subsets $T$ of the remaining pairs.  The degree-$j+2m$ block
is indexed by $(t+m)$-subsets $K$, and its matrix is the transpose of the
inclusion matrix $W_{t,t+m}(v)$.  Each such block occurs with multiplicity
$\binom ns2^s$.

\begin{proof}
Expanding $\xi^m/m!$, terms meeting an occupied pair vanish.  Each empty
$m$-set contributes the basis vector indexed by $T\cup R$, with coefficient
one.  A complete pair has even exterior degree, so moving it past any
other complete pair or singleton introduces no sign.  The singleton set,
its choices, and $T$ uniquely describe every exterior basis vector.
Different such singleton data never mix, proving the integral direct sum.
\end{proof}

For $t\le k\le v-t$, Wilson's diagonal entries for $W_{t,k}(v)$ are
\[
 \binom{k-r}{t-r}\quad\hbox{with multiplicity}\quad
 \binom vr-\binom v{r-1},\qquad 0\le r\le t.
\]
The exact verified primary locator is
\cite[Corollary 3, p.~10]{Ghorbani}, which proves the diagonal form of
\cite{Wilson}.  The original Wilson proof was not opened in this run.
The range condition is necessary; H8's unrestricted version omits it.
For hard Lefschetz, $j=i\le n$ and $m=n-i$, every block has $k=v-t$,
so the condition holds.  A diagonal form is not necessarily a Smith form:
entries $2,3$ normalize to invariant factors $1,6$.

\subsection{SL-02: exact integral hard-Lefschetz consequences}
\label{sl:integral}

For $0\le i\le n$, the map
$L^{(n-i)}:\bigwedge^iM\to\bigwedge^{2n-i}M$ is a
$\mathbb Z$-isomorphism exactly for $i=0,1,n$, with repetitions understood.
For $i=2\le n$, its cokernel is $\mathbb Z/(n-1)$.  Over a commutative
ring $A$, all the \emph{ordinary} hard-Lefschetz maps $L^{n-i}$ are
isomorphisms if and only if $n!$ is a unit in $A$.

\begin{proof}
For $i=0,1$, all blocks have $t=0$, and are $1\times1$ unit matrices.
For $i=n$ the operator is the identity.  If $2\le i<n$, choose
$s=i-2,t=1,v=n-i+2=m+2$.  The block $W_{1,v-1}(v)$ has a diagonal entry
$v-1=m+1>1$, so its cokernel is nonzero.  If $i=2$, this is the only
nonunit block, with $v=n$.  Complementing its column indices gives
$J_n-I_n$, whose Smith form is $(1,\ldots,1,n-1)$.

Necessity of the final assertion follows from $L^n$ on degree zero, which
is multiplication by $n!$.  For sufficiency, the corresponding diagonal
entry of an ordinary map is
\[
 m!\binom{t+m-r}{t-r}=\frac{(t+m-r)!}{(t-r)!}.
\]
It is a product of integers between $1$ and $n$.  If $n!$ is a unit,
each factor is a unit, and the integral diagonal equivalences remain
invertible after base change.
\end{proof}

For a principally polarized complex abelian variety of dimension $g$,
the integral exterior description of cohomology therefore gives cokernel
$\mathbb Z/(g-1)$ for cup product by $\theta^{g-2}/(g-2)!$ on degree
two.  This is a cohomological assertion.  It does not assert algebraicity
of arbitrary divided powers in the Chow ring.

\subsection{SL-03: exact matrices for small powers}
\label{sl:computations}

The reproducible script \texttt{wilson\_checks.py} constructs and checks
16 incidence matrices, aggregates all 27 hard-Lefschetz cases $1\le n\le6$,
and independently constructs the full signed exterior matrices in all 14
cases $n\le4$.  All comparisons pass in exact integer arithmetic.  Smith
forms of diagonal matrices are computed by sorting the valuations separately
at each prime; the direct exterior matrices are checked independently.
The complete results are in the three \texttt{wilson\_*\_checks.tsv}
tables.  Here $1^r$ means $r$ unit invariant factors.

\begin{center}
\begin{tabular}{rrrll}
\toprule
$n$&$i$&$m=n-i$&Smith form of $L^m/m!$&Smith form of $L^m$\\
\midrule
2&0&2&$1$&$2$\\
3&2&1&$1^{14},2$&$1^{14},2$\\
4&2&2&$1^{27},3$&$2^{27},6$\\
4&3&1&$1^{48},2^8$&$1^{48},2^8$\\
5&2&3&$1^{44},4$&$6^{44},24$\\
5&3&2&$1^{110},3^{10}$&$2^{110},6^{10}$\\
5&4&1&$1^{166},2^{43},6$&$1^{166},2^{43},6$\\
6&3&3&$1^{208},4^{12}$&$6^{208},24^{12}$\\
6&4&2&$1^{429},3^{65},6$&$2^{429},6^{65},12$\\
6&5&1&$1^{584},2^{196},6^{12}$&$1^{584},2^{196},6^{12}$\\
\bottomrule
\end{tabular}
\end{center}

For example the degree-two block at $n=4,m=2$ is $J_4-I_4$; the other
24 degree-two basis directions give unit blocks.  Its Smith form is
$1^{27},3$.  These are exact computations in the specified exterior model.
The integer counting theta lifts in step 10 has not thereby been identified
with the power of this operator.

\subsection{SL-04: ordinary contraction and normalized lowering differ}
\label{sl:contraction}

Set $D=\sum_a\iota_{f_a^\vee}\iota_{e_a^\vee}$ and let $H$ act by
$n-j$ on degree $j$.  Then
\[
 [D,L]=H,\qquad[H,L]=-2L,\qquad[H,D]=2D.
\]
For a primitive vector $v$ of degree $i\le n$, meaning $Dv=0$, put
$d=n-i$.  For $0\le r\le d$,
\[
 DL^rv=r(d-r+1)L^{r-1}v,\qquad
 D^rL^rv=r!\frac{d!}{(d-r)!}v.
\]
\begin{proof}
The commutator is checked on the basis $1,e_a,f_a,e_a\wedge f_a$ of
each exterior factor and then summed.  The other two relations follow from
degrees.  For $r=1$, $DLv=Hv=dv$.  If the first formula holds for $r$, then
\[
 DL^{r+1}v=(LD+H)L^rv
 =\bigl(r(d-r+1)+d-2r\bigr)L^rv
 =(r+1)(d-r)L^rv.
\]
Iteration gives the product in the second formula.
\end{proof}
The normalized inverse-string operator satisfies
$\Lambda_N L^rv=L^{r-1}v$; it differs from $D$ on this step by the
factor $r(d-r+1)$.  The scalar $(r!)^2$ in H8 is only the endpoint case
$d=r$.

For $n=2$, put $\Omega=e_1\wedge f_1\wedge e_2\wedge f_2$.  Then
\[
 L^2(1)=2\Omega,\qquad D^2\Omega=2,\qquad D^2L^2(1)=4.
\]
Thus the ordinary-contraction definition in cote 16, scan 014, cannot
simultaneously satisfy the inverse identity on scan 016.  This flags a
normalization, coefficient-ring or revision-layer issue in the supplied
text; it is not a claim about unseen handwriting.  A correct normalized
universal construction is given in SL-08 below.

\subsection{SL-05: Arthur blocks do not determine integral indices}
\label{sl:arthur-lattice}

An Arthur block $\mu[t]$ has the form
$\mu\boxtimes\operatorname{Sym}^{t-1}(\mathbb C^2)$.  Its $\mathrm{SL}_2$
weights are $t-1,t-3,\ldots,1-t$, and its associated unramified exponents
are their halves \cite[\S2.3, equation (2.3.1)]{Mok}.  The block is not
itself an integral motive with cup-product correspondences.

\paragraph{Integral counterexample, including $m=1,2,3$.}
In $V_d=\operatorname{Sym}^d(\mathbb Q^2)$, consider the monomial and
binomial lattices
\[
 \Gamma_0=\bigoplus_{r=0}^d\mathbb Zx^{d-r}y^r,\qquad
 \Gamma_1=\bigoplus_{r=0}^d
             \mathbb Z\binom drx^{d-r}y^r.
\]
Both are stable under $\mathrm{SL}_{2,\mathbb Z}$ and give the identical
rational representation.  For $d=m+1$ the divided operator
$e^{(m)}=(x\partial_y)^m/m!$ from the line of $y^d$ to the line of
$x^my$ has Smith invariant $m+1$ on $\Gamma_0$ and $1$ on $\Gamma_1$.

\begin{proof}
The monomial lattice is stable by polynomial substitution.  For the second
lattice, the tensors invariant under permutations in $(\mathbb Z^2)^{\otimes d}$
have the orbit sums with $r$ copies of $y$ as a basis.  Symmetric
multiplication sends these to $\binom drx^{d-r}y^r$ and is equivariant.
Its image is therefore a stable lattice.  Direct differentiation gives
$e^my^{m+1}=(m+1)!x^my$.  The target basis vector in $\Gamma_1$ is
$(m+1)x^my$, so the ordinary coefficients are $(m+1)!$ and $m!$;
division by $m!$ gives $m+1$ and $1$.
\end{proof}

\begin{center}
\begin{tabular}{rcll}
\toprule
$m$&Rational representation&Ordinary indices $(\Gamma_0,\Gamma_1)$
&Divided indices $(\Gamma_0,\Gamma_1)$\\
\midrule
1&$\operatorname{Sym}^2$&$(2,1)$&$(2,1)$\\
2&$\operatorname{Sym}^3$&$(6,2)$&$(3,1)$\\
3&$\operatorname{Sym}^4$&$(24,6)$&$(4,1)$\\
\bottomrule
\end{tabular}
\end{center}
The script also constructs these raising matrices and verifies the table;
the output is \texttt{sl2\_lattice\_checks.tsv}.  This is a fully specified
counterexample to a lattice-free inference.  It is not a numerical claim
about the actual theta lift.

A theta chain changes the automorphic group and adjoins blocks, whereas
cup product moves inside a fixed graded realization.  Identifying the
number of lifts with an Arthur-block length or a Lefschetz exponent requires
additional data.  An index of the associative order generated by operators
also needs a specified ambient integral order; it is not determined by the
Smith form of one operator.

\subsection{SL-06: the actual local divisibility statement and its parameters}
\label{sl:actual-code}

The retrieved \path{Formalism/LocalToGlobal/Divisibility/Inputs.lean}
defines \texttt{ColChain} using a length \texttt{mm}, varying dimensions
$c(j)$, rectangular matrices $K_j$ of size $c(j+1)\times c(j)$, and a
final period functional.  Its field \texttt{fact} identifies the original
period with that functional on the product of the matrices.  It includes
no Lefschetz grading, lowering operator or $\mathfrak{sl}_2$ relation.

Here is the exact scalar content of the inspected local estimate.  Let
$w\mid\ell$ be a place of a number field $K$, with ramification index
$e_w$, normalized by
$|\iota(a)|_\ell=|\ell|_\ell^{v_w(a)/e_w}$.  Suppose $J$ is a finite
nonempty set with a distinguished $j_1$, and a period has a finite expansion
\[
 \iota(P)=\sum_\pi x_\pi y_\pi r_\pi\gamma_{P,\pi}^{-1}
        \prod_{j\in J}
          (\gamma_{\pi,j}^{-1}\varepsilon_{\pi,j}a_{\pi,j}).
\]
Assume that the norms of $x_\pi,y_\pi,r_\pi,\varepsilon_{\pi,j}$ are
at most one; $D$ is a nonzero natural number; all denominators
$\gamma_{\pi,j},\gamma_{P,\pi}$ divide $D$;
$\gamma_{\pi,j_1}=1$; and $|a_{\pi,j}|\le|\ell|^{s_j}$.  Then
\[
 P=0\quad\hbox{or}\quad
 v_w(P)\ge e_w\sum_{j\in J}(s_j-v_\ell(D)).
\]
\begin{proof}
Each inverse denominator costs at most $|\ell|^{-v_\ell(D)}$ in norm.
There are at most $|J|-1$ such costs in the product because the distinguished
denominator is one, and at most one further cost from $\gamma_{P,\pi}$.
Every summand is consequently bounded by
$|\ell|^{\sum_js_j-|J|v_\ell(D)}$.  The ultrametric inequality gives
the same bound for the sum.  If $P\ne0$, substituting the valuation
normalization and using $0<|\ell|<1$ gives the claimed inequality.
\end{proof}

The exact locators within the same directory are
\path{JetBounds/IsExpansion.lean}, declarations
\texttt{PathData}, \texttt{PathData.val}, \texttt{StabDvd},
\texttt{IsExpansion}; \path{JetBounds/Statement.lean}, declarations
\texttt{PathBound}, \texttt{LocalValuationBoundAt},
\texttt{PlaceInputAtEntry}; and
\path{JetBounds/JetBoundsProofs/ValuationOfEntryBounds.lean}, theorem
\texttt{valuation\_of\_entry\_bounds}.  That complete proof and the
complete proofs of \texttt{path\_height\_bound},
\texttt{embedding\_lower\_bound} and
\texttt{norm\_inv\_natCast\_le\_of\_dvd} in
\texttt{NormInvNatCastLeOfDvd.lean} were read.  They implement the
argument above.  The upstream entry estimates and full global application
were not re-proved, and no kernel execution is claimed.

For $|J|=1,2,3$, the bounds are respectively
\[
 e_w(s_1-v_\ell D),\quad
 e_w(s_1+s_2-2v_\ell D),\quad
 e_w(s_1+s_2+s_3-3v_\ell D).
\]
These specialize the actual statement while displaying its parameters.
They are lower bounds, not exact valuations or Smith factors.  Integer-valued
$v_w$ permits replacing the right side by its ceiling.

\paragraph{Actual small-parameter data.}
The arithmetic family in
\path{Formalism/LocalRingCardinality/Params.lean}, structure
\texttt{Ranks}, assumes $m$ even and $m\ge2$, as well as
$m+3\le k$, $k+m+4\le t$, $3m<t$, and
$m\le2\operatorname{ord}\le k$.  Thus the $m=1,3$ examples above are
universal algebraic models, not instances of this parameter family.

The concrete \texttt{numEx} in \texttt{NumbersLemmas.lean} has
$m=2,k=8,t=15,nV=24$, host $q=11,d=7$, and distinguished primes
$3697,7393$.  Its ranks are $N(j)=8+16j$, so the two steps have ranks
$8\to24\to40$.  The exact underlying declarations are
\texttt{R0} in \path{SpanPowerMembership/ExampleRanksRankComplement.lean}
and \texttt{hcV}, \texttt{hcB}, \texttt{hcE} in
\path{SpanPowerMembership/ExampleWeightMiddleEntries.lean}, both under
\path{Formalism/LocalRingCardinality}.  They give
\[
 v=(24,20,16,12,8,4,0),\quad b=(91,109),\quad e=(100,118).
\]
The host index set has four elements, so the real-place index set has
eight.  At each index, the rows in \texttt{Numbers.lean} are
\[
\begin{array}{c|l|l}
j&a^+_j&a^-_j\\ \hline
1&(109,106,103,100,97,94,91,1)&(1)\\
2&(126,123,120,117,114,111,108,18,1)&()
\end{array}
\]
and therefore
\[
 (D^+_1,D^+_2)=(5608,6704),\quad
 (D^-_1,D^-_2)=(8,0),
\]
\[
 (F_{G,1},F_{G,2})=(-2808,-3352),\qquad
 (F_{\mathrm{ram},1},F_{\mathrm{ram},2})=(-1400,-1652).
\]
These follow by exact rational substitution, with
$F_{G,j}=-(D^+_j+D^-_j)/2$ and
$F_{\mathrm{ram},j}=-(7/2)\cdot4e_j$.
The first ramified coefficient and the first zero coefficients at $2,11$
also agree with the explicit source declarations
\texttt{ex\_FRamD}, \texttt{ex\_FjP\_d},
\texttt{ex\_FjP\_2}, \texttt{ex\_FjP\_q}.

\paragraph{Exact evaluation of the orbit optimization.}
For constant inputs $D_s=A\ge0$, $p\ge2$ and an orbit length
$d_0\ge2$, the linear program \texttt{lpFeasible} in
\texttt{DataParams\textsubscript{2}.lean} has optimum
$d_0A/(p+1)$, and
\[
 \operatorname{lpF}(p,d_0,A)=\frac{d_0A}{p^2-1}.
\]
\begin{proof}
The variables $x_{s,i}\ge0$, indexed by $s\in\mathbb Z/d_0\mathbb Z$,
vanish unless $i<d_0$ is odd.  The constraints are
\[
 \sum_{i<d_0}p^i x_{s,i}+\sum_{i<d_0}x_{s-i,i}\le A.
\]
Summing over $s$ and translating the second sum gives
\[
 (p+1)\sum_{s,i}x_{s,i}
 \le\sum_{s,i}(p^i+1)x_{s,i}\le d_0A,
\]
since every supported $i$ is at least one.  The point
$x_{s,1}=A/(p+1)$, all other variables zero, satisfies every constraint
with equality and attains this bound.  Substituting the optimum in
$\operatorname{lpF}=p d_0A/(p^2-1)-\sup\operatorname{lpFeasible}$
gives the formula.
\end{proof}
Here $A=7e_j$.  If the inert class has zero Frobenius element, the
exceptional branch of \texttt{FInert} gives the same result directly.
In the remaining inert or signal cases the orbit length is at least two,
and the displayed lemma applies.  Consequently, at any such prime,
\[
 F_{\mathrm{inert},j}(p)=\frac{56e_j}{p^2-1}.
\]

For \texttt{numEx}, the cutoff is
$X''=8\exp(2^{1/300})$, with $21<X''<22$.
For a rational certificate, $(101/100)^{300}>2$ and the Taylor sum
through degree $12$ for $\exp(101/100)$, with the geometric bound on its
remaining tail, is less than $11/4$; also $e>8/3$.
The prime classification is the one in \texttt{HostCls.lean}:
outside $2,7,11$, split primes have $(-7/p)=1$; the others are signal
or inert according as $(-11/p)=1$ or $-1$.
This gives the entire nonzero contribution table for this numerical datum:
\begin{center}
\begin{tabular}{rcrrr}
\toprule
$p$&Class&$F_{1,p}$&$F_{2,p}$&$\sum_jF_{j,p}$\\
\midrule
2&ramified at $2$&0&0&0\\
3&signal&700&826&1526\\
5&signal&$700/3$&$826/3$&$1526/3$\\
7&ramified at $d$&$-1400$&$-1652$&$-3052$\\
11&ramified at $q$&0&0&0\\
13&inert&$100/3$&$118/3$&$218/3$\\
17&inert&$175/9$&$413/18$&$763/18$\\
19&inert&$140/9$&$826/45$&$1526/45$\\
3697,7393&distinguished&$-2808$&$-3352$&$-6160$\\
\bottomrule
\end{tabular}
\end{center}
For each of the last two primes the row applies separately.  In particular,
the actual definition of \texttt{Fo} evaluates to
\[
 12208\sum_{p\in\{3,5,13,17,19\}}\frac{\log p}{p^2-1}
 -3052\log7-6160(\log3697+\log7393).
\]
Its decimal approximation is $-108518.601323$; this is the $F$
contribution alone.  The exact evaluation, the cutoff certificate and
twenty primal/upper-bound certificates for the relevant orbit programs
are reproduced by \texttt{step10\_numeric\_checks.py}.  The parameter
and prime tables are saved separately with hashes of the inspected source
files.  No Lean kernel execution is claimed.

\paragraph{The exact remaining hypotheses.}
The small tuple \texttt{numEx} is a test datum, and cannot satisfy
\texttt{NumFacts} in \texttt{NumbersLink.lean}: its field \texttt{n2}
would require $N(m)^2\le b$, whereas
\[
 N(m)^2=40^2=1600,\qquad b=\exp(2^{1/100})<e^2<9.
\]
This proves incompatibility of that particular tuple with these hypotheses;
it does not refute the conditional theorem.  The same file explicitly
provides the large-scale repair \texttt{NumInst.numBig}, with
$M=200+|\texttt{numEx.logOm}|$ and $x=\exp(100M)$, and constructs
\texttt{numFactsBig : numBig.NumFacts wd0 1 1 5 5 1 53 4}.
Its $v,b,e$ are unchanged; only one $u$ entry and the scale change.
Thus the distinguished and ramified coefficients above and the formula
$56e_j/(p^2-1)$ also hold for this larger example, with further primes
admitted by its larger cutoff.  This is a check of the exact declaration
and its parameters, not a full check of all supporting \texttt{numFactsBig}
proofs.

Neither a numerical datum nor \texttt{NumFacts} selects a nonzero period
and its two integral lattices.  The exact record-level definition in
\texttt{DataParams\textsubscript{2}.lean} is
\[
 \operatorname{sg}(u,w)=\sum_{j=1}^m
   \bigl(\operatorname{gainD}(j,w)+F(j,w)-E(u,j,w)\bigr),
\]
where \texttt{gainD} suppresses the gain at primes in \texttt{drop}.
The global gain term also subtracts \texttt{tau}; it is not part of this
local definition.  One still needs the record linking its $F$ to the
numeric $F_{j,p}$, the gain and error data, the period expansion, entry
estimates, integral row and column lattices, stabilizer denominator, and
the inequality relating those inputs to the requested local threshold.
For comparison, the companion \texttt{ExampleWeightMiddleEntries.lean}
also constructs an \texttt{emptyColumnChain} and
\texttt{trivialPeriodMap} whose period is identically zero.  These examples
are not a supplied nonzero theta-lift matrix with a specified Smith form.

At the level of matrix-chain data and factorization, there is no invariant
depending only on length.  Fix $m\ge1$, set every dimension equal to one,
take all frames and all matrices equal to one except $K_0=[\ell^s]$, and
take the final functional to be multiplication.  The period is
$\ell^sxy$, satisfying the factorization and taking valuation $s$ at
$x=y=1$.  Every $s\ge0$ is possible at fixed length.  This refutes a
length-only conclusion from that bare matrix structure, not from the
additional automorphic hypotheses of a particular lift.

\paragraph{The actual divided-power action.}
In \path{Formalism/GradedWeylDivisibility/StatementA.lean},
\texttt{rowDiv R a b k f} is the coefficient of $T^k$ after the row
transvection $X_{b,s}\mapsto X_{b,s}+TX_{a,s}$.  The Weyl lattice is
generated by words in these operators on a highest-weight polynomial;
the dual Weyl lattice is defined by column homogeneity and unipotent
invariance.  \texttt{StatementB.lean} uses a factorial-weighted Fock
pairing, and \texttt{StatementD.lean}, \texttt{StatementE.lean}
treat transformed integral lattices and duality.  These are concrete
integral general-linear-group modules.  No identification with the
symplectic exterior operator $L$ is present in the inspected definitions.

An equality with Wilson's factors therefore needs a map to a fixed graded
geometric realization, equality of the actual operators with specified
divided powers, and preservation of the specified integral lattices and
denominators.  The numerical input not yet supplied is a comparison gate,
not itself a mathematical counterexample.  A scoped search of the three
divisibility folders, \texttt{PositivityTheory} and
\path{Formalism/HermitianCurvature} found no named arithmetic Hodge-index
comparison; this is not a nonexistence assertion about the entire archive.

\subsection{SL-07: a Frobenius normalization obstruction}
\label{sl:frobenius}

Let $X/\mathbb F_q$ be smooth projective of positive dimension, and let
$\eta$ be an ample hyperplane class defined over $\mathbb F_q$.  Use
cohomological geometric Frobenius $F$ on untwisted $\ell$-adic cohomology,
so $F\eta=q\eta$.  Then for $Lx=\eta\cup x$,
\[
 FL=qLF.
\]
\begin{proof}
Pullback by the $q$-power Frobenius sends $\mathcal O(1)$ to
$\mathcal O(q)$.  Chern-class functoriality gives $F\eta=q\eta$,
and cup-product functoriality proves the identity.  Equivalently, the class
is Galois-invariant in $H^2(\overline X,\mathbb Q_\ell(1))$.
\end{proof}
The conventions and Chern-class facts were checked in
\cite[Theorem 23.3 and \S27, especially Proposition 27.3]{MilneEtale}.

In contrast, the two factors of a local parameter
$\psi:L_F\times\mathrm{SL}_2(\mathbb C)\to{}^LG$ commute.  Thus the
raising element commutes with $\psi(w,1)$.  A linear identification of
that pair with the untwisted pair $(L,F)$ would imply
$(q-1)LF=0$.  Since $F$ is invertible and coefficients have characteristic
zero, it would force $L=0$.  Already $\mathbb P^1$ gives
$F=\operatorname{diag}(1,q)$ and a nonzero map $1\mapsto\eta$.

This only obstructs the unnormalized identification.  On
$\bigoplus_rH^{2r}(\overline X,\mathbb Q_\ell(r))$, and separately on
$\bigoplus_rH^{2r+1}(\overline X,\mathbb Q_\ell(r))$, hyperplane
multiplication commutes with Frobenius.  These integer Tate twists are
also integral over $\mathbb Z_\ell$ when $\ell\ne\operatorname{char}\mathbb F_q$.
One must not call every such normalization nonintegral.  A uniform
weight-zero normalization can additionally require half twists.
Furthermore the associated Langlands parameter is
\[
 \varphi_\psi(w)=\psi\left(w,
 \begin{pmatrix}|w|^{1/2}&0\\0&|w|^{-1/2}\end{pmatrix}\right),
\]
not $\psi(w,1)$, and need not commute with the raising element
\cite[\S2.2, equations (2.2.9)--(2.2.10)]{Mok}.  Proper normalization
removes this particular obstruction, but does not provide the missing
integral geometric comparison.

\subsection{SL-08: the normalized universal ring, with complete proof}
\label{sl:universal}

Define
\[
 \Phi_n^{\mathrm{str}}=\prod_{i=0}^n M_{n-i+1}(\mathbb Z).
\]
Its $i$th factor acts on a string $v_{i,0},\ldots,v_{i,n-i}$ in degrees
$i,i+2,\ldots,2n-i$.  Let $L$ move one step up and $\Lambda_N$ one
step down, both zero outside the string.  Write $E_i^{\beta,\alpha}$
for its matrix units, of degree $2(\beta-\alpha)$, and let $\pi_j$
be the sum of diagonal units in degree $j$.  Then
\[
 E_{i'}^{\beta',\alpha'}E_i^{\beta,\alpha}
   =\delta_{i',i}\delta_{\alpha',\beta}E_i^{\beta',\alpha}.
\]
The ring is generated over $\mathbb Z$ by $L,\Lambda_N$ and is free of
rank
\[
 \sum_{i=0}^n(n-i+1)^2=\frac{(n+1)(n+2)(2n+3)}6.
\]

\begin{proof}
Set $B=1-L\Lambda_N$, the projection to all string bottoms.  For
$0\le d\le n$,
\[
 Q_d=B(\Lambda_N^dL^d-\Lambda_N^{d+1}L^{d+1})B
\]
projects to the bottom of the string of length $d+1$.  Indeed the two
terms on a bottom vector are the indicators for lengths at least $d+1$
and at least $d+2$.  Every matrix unit on this string is
$L^\beta Q_d\Lambda_N^\alpha$.  This proves generation, while the
product of matrix rings gives the asserted free rank.
\end{proof}

The correct string-index range is $0\le\alpha,\beta\le n-i$.
The range $2n-i$ in scan 015 is incompatible with its degrees and its
claimed basis.  The corrected rank is five for $n=1$, agreeing with
scan 017.

\paragraph{Universal property.}
Let $A$ be a commutative ring and $M=\bigoplus_{j=0}^{2n}M^j$ a graded
$A$-module, with a degree-two $A$-linear operator $L$.  Suppose the ordinary
maps
\[
 L^{n-i}:M^i\longrightarrow M^{2n-i}
\]
are isomorphisms for $0\le i\le n$.  Then there is a unique graded
$A$-linear representation of $A\otimes\Phi_n^{\mathrm{str}}$ carrying
$L$ and every degree projector to the prescribed operators.  Its lowering
operator is the normalized inverse-string operator.  Conversely, such a
representation implies these hard-Lefschetz isomorphisms.  For a
unital graded $A$-subalgebra $\mathcal E\subset\operatorname{End}_A(M)$
containing $L$, existence of the
representation into $\mathcal E$ is equivalent to the normalized lowering
operator and all degree projectors belonging to $\mathcal E$.

\begin{proof}
For $i\le n$, put
\[
 P^i=\ker(L^{n-i+1}:M^i\to M^{2n-i+2}),
\]
with out-of-range groups zero.  If $i\ge2$, the map
\[
 x\longmapsto
 (L^{n-i+2}:M^{i-2}\to M^{2n-i+2})^{-1}L^{n-i+1}x
\]
is a left inverse of $L:M^{i-2}\to M^i$, with kernel $P^i$.
Therefore $M^i=P^i\oplus LM^{i-2}$; for $i=0,1$ one has $P^i=M^i$.
Induction proves the decomposition in degrees at most $n$.  Transporting
the decomposition of degree $2n-j$ through $L^{j-n}$ proves it in degree
$j>n$.  In total,
\[
 M=\bigoplus_{i=0}^n\bigoplus_{r=0}^{n-i}L^rP^i.
\]
Each $L^r$ is injective on $P^i$ because $L^{n-i}$ is injective there.
Thus matrix units act unambiguously by $L^\alpha v\mapsto L^\beta v$
within each primitive string.  This constructs the representation and
the lowering operator, and immediately gives the converse.

For uniqueness, a representation with the prescribed degree projectors
decomposes $M$ into its matrix-unit summands.  Its bottom summand in the
$i$th factor is exactly $P^i$: $L^{n-i+1}$ kills that summand and is
injective on every other summand in degree $i$.  The remaining summands
are consequently $L^rP^i$, and $L$ forces every matrix unit between them.
The representation is unique.  Since the matrix units are generated by
$L,\Lambda_N$, the subalgebra assertion follows.
\end{proof}

The degree condition is essential.  For $n=1$, conjugation by $1+tL$ in
the $M_2(\mathbb Z)$ factor fixes $L$ but sends
\[
 \Lambda_N\longmapsto
 \Lambda_N+t(E_{11}-E_{00})-t^2L.
\]
Different integers $t$ give different ungraded maps with the same $L$;
they do not preserve the degree projectors.

The symplectic exterior lattice with ordinary $L$ does not furnish this
unit-step model over $\mathbb Z$ in general.  By SL-02, all its ordinary
hard-Lefschetz maps become isomorphisms after inverting $n!$.  Already for
$n=2$, $L^2(1)=2\Omega$ forces $\Lambda_N\Omega=\xi/2$, which is not
integral in the exterior lattice.  After base change to
$\mathbb Z[1/n!]$, the universal property applies.  This repairs the
normalized calculation; it does not identify this ring with the distinct
integral contraction order $\mathbb Z[L,D]$ in an unspecified ambient
matrix order.

\paragraph{Verification boundary.}
The new elementary arguments above were read in full by the authoring
agent.  AR-01 was independently reconstructed and its entire proof checked
by the openness agent; SL-08 was independently re-proved by the sign-geometry
agent.  Literature statement checks, complete proof checks, exact integer
computations and source-level Lean inspection are recorded separately in
the accompanying event ledger.  No full Lean compilation or proof of the
entire imported dependency graph is claimed.


\section{A specified integral operator order: exact indices in small rank}
\label{sec:order}

The order-index question can be answered once the lattice, lowering
operator, and ambient order are fixed.  This section supplies one
natural precise model and a complete elementary computation.  It
separates two meanings of lowering that have the same full rational
associative envelope but different integral orders.  No identification
with the integral lattice of a theta-lift chain is assumed.

\begin{theorem}[SL-09: indices in the ordinary-power string model]
\label{sl:order}
For \(n\ge1\), set
\[
\mathcal A_n=\prod_{i=0}^n
  \operatorname{Mat}_{n-i+1}(\mathbf Z).
\]
In factor \(i\), put \(d=n-i\), and choose basis
\(v_{i,0},\ldots,v_{i,d}\) in degrees
\(i,i+2,\ldots,2n-i\).  Define the integral operators
\[
L v_{i,r}=v_{i,r+1},\qquad
D v_{i,r}=r(d-r+1)v_{i,r-1},
\]
with zero outside the indicated ranges.  Let \(\pi_j\) be the
projector to degree \(j\), and let
\[
\mathcal O_n=\mathbf Z[L,D,\pi_0,\ldots,\pi_{2n}]
\subseteq\mathcal A_n
\]
be the unital associative subring.  Then
\[
[\mathcal A_1:\mathcal O_1]=1,\qquad
[\mathcal A_2:\mathcal O_2]=2^5=32,\qquad
[\mathcal A_3:\mathcal O_3]=2^{13}3^{10}=483729408.
\]
If \(D\) is replaced by the normalized unit lowering
\(\Lambda_N v_{i,r}=v_{i,r-1}\), the generated order is
\(\mathcal A_n\) for every \(n\), and its index is \(1\).
\end{theorem}

We give a direct description of the two nontrivial orders.  Matrix
indices within a string start at \(0\); \(E_{rs}\) means row \(r\),
column \(s\).  Thus \(L\) has its nonzero entries just below the
diagonal, and \(D\) has its nonzero entries just above it.

\subsection{An elementary divisibility order}

For a string of length \(d+1\), put
\[
c_r=r(d-r+1)\quad(1\le r\le d),\qquad
u_0=1,\quad u_r=\prod_{a=1}^r c_a.
\]
Define
\[
T_d=\left\{(a_{rs})\in\operatorname{Mat}_{d+1}(\mathbf Z):
       \frac{u_s}{u_r}\mid a_{rs}\text{ whenever }r<s\right\}.
\]
There is no divisibility condition on or below the diagonal.

\begin{lemma}
\label{sl:triangular-order}
The set \(T_d\) is a subring.  It is generated by \(L,D\) and all
the individual diagonal projectors of this single string.  Its
index in the full matrix ring is
\[
\prod_{0\le r<s\le d}\frac{u_s}{u_r}.
\]
\end{lemma}

\begin{proof}
For \(r<s\), consider a term \(a_{rk}b_{ks}\) in the
\((r,s)\)-entry of a product.  If \(r\le k\le s\), the two required
divisibilities supply the product
\((u_k/u_r)(u_s/u_k)=u_s/u_r\).  If \(k<r\), the second factor is
divisible by \(u_s/u_k\), itself a multiple of \(u_s/u_r\).
If \(k>s\), the first factor is divisible by \(u_k/u_r\), again
a multiple of \(u_s/u_r\).  Thus every summand meets the required
condition; \(T_d\) is a subring.

With individual diagonal projectors available, powers of \(L\)
give \(E_{rs}\) for \(r\ge s\), and powers of \(D\) give
\((u_s/u_r)E_{rs}\) for \(r<s\).  These form a
\(\mathbf Z\)-basis of \(T_d\).  The displayed product is therefore
the index of this diagonal coordinate sublattice in the full
matrix ring.
\end{proof}

The degree projectors in the theorem are not generally the
individual diagonal projectors: different strings can contribute
to the same degree.  The resulting congruence between factors
accounts for the remaining index.

\subsection{The cases n=1 and n=2}

The sums of the even and odd degree projectors are central
idempotents in the ambient ring and belong to \(\mathcal O_n\).
Consequently \(\mathcal O_n\) splits into its two parity factors.

For \(n=1\), the length-two string occupies degrees \(0,2\);
the length-one string occupies degree \(1\).  These factors are
separated by degree projectors.  The lowering coefficient on
the length-two string is \(1\).  Hence both factors are full
matrix rings, and the index is \(1\).

For \(n=2\), the odd factor is a single length-two string and
is again the full \(\operatorname{Mat}_2(\mathbf Z)\).
The even factor consists of a length-three string in degrees
\(0,2,4\) and a length-one string in degree \(2\).
For the first string,
\[
(c_1,c_2)=(2,2),\qquad (u_0,u_1,u_2)=(1,2,4).
\]
Thus \(T_2\) imposes
\[
2\mid a_{01},\qquad 2\mid a_{12},\qquad 4\mid a_{02},
\]
and has index \(2\cdot2\cdot4=2^4\) in
\(\operatorname{Mat}_3(\mathbf Z)\).

The map
\[
\rho_2:T_2\longrightarrow\mathbf F_2,\qquad
A\longmapsto a_{11}\bmod2
\]
is a surjective ring homomorphism.  Indeed in
\((AB)_{11}\), the two terms through indices \(0\) and \(2\)
are even by the displayed divisibilities.  The even order is
exactly
\[
T_2\times_{\mathbf F_2}\mathbf Z
=\{(A,z):a_{11}\equiv z\pmod2\}.
\tag{SL-09a}
\]

To prove equality, first observe that all of \(L,D,\pi_j\)
satisfy the description, so the generated ring is contained in
the right side.  Conversely, the endpoint projectors give
\((E_{00},0)\) and \((E_{22},0)\).  Multiplying \(L\) and \(D\)
by these and the middle degree projector gives
\[
(E_{10},0),\quad(E_{21},0),\quad
(2E_{01},0),\quad(2E_{12},0).
\]
Their products give \((E_{20},0)\), \((4E_{02},0)\), and
\((2E_{11},0)\).  The remaining degree projector is
\((E_{11},1)\).  These elements span every pair satisfying
\((\mathrm{SL\text{-}09a})\): all the off-diagonal coordinates
have their minimal required divisibility, and every pair of
middle diagonal entries with the same residue modulo \(2\)
is generated by \((E_{11},1)\) and \((2E_{11},0)\).
This proves equality.

Equality of the two residues is one independent index-two
condition on \(T_2\times\mathbf Z\).  Therefore the even factor
has index \(2^4\cdot2=2^5\), and so does \(\mathcal O_2\).

\subsection{The case n=3}

The odd factor is the length-three string in degrees \(1,3,5\)
together with the length-one string in degree \(3\).  It is
exactly the order just computed, with index \(2^5\).

The even factor consists of the length-four string in degrees
\(0,2,4,6\) and the length-two string in degrees \(2,4\).
Use \(E_{rs}\) for the first factor and \(e_{rs}\) for the second.
For the first factor,
\[
(c_1,c_2,c_3)=(3,4,3),\qquad
(u_0,u_1,u_2,u_3)=(1,3,12,36).
\]
The order \(T_3\) has index
\[
3\cdot12\cdot36\cdot4\cdot12\cdot3
=2^8 3^6
\tag{SL-09b}
\]
in \(\operatorname{Mat}_4(\mathbf Z)\).

Let
\[
\rho_3:T_3\longrightarrow\operatorname{Mat}_2(\mathbf F_3)
\]
be reduction modulo \(3\) of the submatrix on rows and columns
\(1,2\).  This is a ring homomorphism.  In a product of two
matrices, contributions to this submatrix through index \(0\)
contain either \(b_{01}\) or \(b_{02}\), both divisible by \(3\).
Contributions through index \(3\) contain either \(a_{13}\) or
\(a_{23}\), likewise divisible by \(3\).  The remaining terms
are precisely the product of the two middle submatrices.
The map is surjective: the middle diagonal entries and the
lower off-diagonal entry are arbitrary, and the upper
off-diagonal entry is a multiple of \(4\), which is a unit
modulo \(3\).

We claim that the even order is
\[
T_3\times_{\operatorname{Mat}_2(\mathbf F_3)}
                 \operatorname{Mat}_2(\mathbf Z).
\tag{SL-09c}
\]
Here the second map is ordinary reduction modulo \(3\).
Every generator satisfies this condition.  For the middle
link, the big-string lowering coefficient is \(4\), while
the small-string coefficient is \(1\); these coincide
modulo \(3\).  The raising coefficients are both \(1\).
The two middle degree projectors have the same matching
diagonal entries.  Endpoint terms disappear under \(\rho_3\).
This proves one containment.

For the reverse containment, endpoint projections and powers
of \(L\) supply the following isolated big-string units:
\[
(E_{00},0),\ (E_{33},0),\
(E_{10},0),\ (E_{20},0),\ (E_{30},0),\
(E_{31},0),\ (E_{32},0).
\]
For example,
\(\pi_6L^2\pi_2=(E_{31},0)\); the small string has no
degree-six component.  Powers of \(D\) similarly supply
\[
(3E_{01},0),\ (12E_{02},0),\ (36E_{03},0),\
(12E_{13},0),\ (3E_{23},0).
\]
These are all coordinates outside the middle \(2\)-by-\(2\)
block, with exactly the minimal divisibilities required
by \(T_3\).

Their products also supply the four middle kernel elements
\[
(3E_{11},0),\quad(3E_{22},0),\quad
(3E_{21},0),\quad(12E_{12},0).
\]
Explicit products, ignoring the zero second components, are
\[
E_{10}(3E_{01}),\qquad
(3E_{23})E_{32},\qquad
E_{20}(3E_{01}),\qquad
E_{10}(12E_{02}),
\]
respectively.  Finally the degree projectors and the middle
raising and lowering links provide four coupled lifts:
\[
(E_{11},e_{00}),\quad (E_{22},e_{11}),\quad
(E_{21},e_{10}),\quad (4E_{12},e_{01}).
\]
Given a pair in \((\mathrm{SL\text{-}09c})\), subtract a
linear combination of these four coupled lifts to kill its
entire small-string component.  The remaining middle
big-string entries are divisible by \(3,3,3,12\), respectively:
the upper entry is already divisible by \(4\), and its
residue modulo \(3\) now vanishes.  They are therefore a
linear combination of the four kernel elements just produced.
The isolated outer coordinates handle the rest.  This proves
\((\mathrm{SL\text{-}09c})\).

Since both residue maps are onto
\(\operatorname{Mat}_2(\mathbf F_3)\), the fibre product has
additive index
\[
\left|\operatorname{Mat}_2(\mathbf F_3)\right|=3^4
\]
in \(T_3\times\operatorname{Mat}_2(\mathbf Z)\).
Combining this with \((\mathrm{SL\text{-}09b})\) and the
odd factor yields
\[
[\mathcal A_3:\mathcal O_3]
=2^5\cdot(2^8 3^6)\cdot3^4
=2^{13}3^{10}.
\]
This completes the non-normalized part of
Theorem~\ref{sl:order}.

\subsection{Normalized lowering and the scope of H8(iv)}

For unit lowering \(\Lambda_N\), the full matrix-unit
generation was proved in SL-08, Section~\ref{sl:universal}.
For completeness, put \(B=1-L\Lambda_N\), which projects
to the bottom of each string.  Then
\[
Q_d
=B\bigl(\Lambda_N^dL^d-\Lambda_N^{d+1}L^{d+1}\bigr)B
\]
is the bottom projector of the unique string of length
\(d+1\).  Its matrix units are
\[
L^\beta Q_d\Lambda_N^\alpha,\qquad
0\le\alpha,\beta\le d.
\]
They generate the full product \(\mathcal A_n\) integrally,
even without including degree projectors as additional
generators.  This proves the final assertion of
Theorem~\ref{sl:order}.

Consequently the index of an integral operator order is
not an invariant of the full rational associative envelope
or the raising operator alone.  In the same specified
ambient order, with the same \(L\) and the same degree
projectors, the two lowering conventions give indices
\(32\) and \(1\) already for \(n=2\).

This refutes the automatic, normalization-independent
inference asserted in H8(iv), that the operator-order
index is supplied by the same binomials as the individual
exterior-algebra maps.  In particular, all divided-power
hard-Lefschetz maps on the symplectic exterior algebra for
\(n=2\) have index \(1\), whereas the explicitly specified
ordinary-power string order here has index \(32\).
These are different lattice constructions, and their
indices cannot be identified without a map respecting
the lattices and the exact operators.

The phrase \emph{with or without divided powers} also needs
a new integrality specification in this ambient order:
on a unit-step string of length at least three,
\(L^2/2!\) has a matrix entry \(1/2\), so it does not
belong to \(\mathcal A_n\).  The theorem does not assign
an index in \(\mathcal A_n\) to that noncontained ring.
It also makes no assertion that \(\mathcal O_n\) is the
image of the integral symplectic exterior lattice or the
integral lattice of an Arthur parameter.  It settles this
particular precise order model and exhibits the
normalization dependence that a further comparison must
resolve.

\subsection{Exact independent computation}

\begin{center}
\begin{tabular}{rrrr}
\toprule
\(n\)&\(\operatorname{rank}_{\mathbf Z}\mathcal A_n\)&
\([\,\mathcal A_n:\mathcal O_n\,]\)&
Index for \(\Lambda_N\)\\
\midrule
1&5&1&1\\
2&14&32&1\\
3&30&483729408&1\\
\bottomrule
\end{tabular}
\end{center}

The companion program
\texttt{normalized\_order\_indices.py} independently computes
these indices by exact column Hermite normal form.  It
vectorizes each ambient matrix unit, begins with the
\(\mathbf Z\)-span of \(1\), and repeatedly replaces a
basis matrix \(B\) by the column HNF of
\[
\bigl[B,\ R_LB,\ R_DB,\ R_{\pi_0}B,\ldots,
                         R_{\pi_{2n}}B\bigr],
\]
where \(R_g\) is right multiplication by \(g\).
At each stage this lattice is precisely the span of the
words reached by adjoining one more generator, together
with all previously reached words.  Equality with the
previous HNF therefore proves stability under every
generator, and hence equality with the generated
unital order.  The stable matrices have full rank;
their determinant absolute values are the indices.

For \(n=1,2,3\), stabilization was detected on iterations
\(2,3,4\), respectively, for both lowering conventions.
The script separately verifies, with exact rational
arithmetic, that each final lattice contains \(1\) and
all generators and that \(B^{-1}R_gB\) has integer
entries for every generator.  No floating point
determinant or probabilistic lattice test is used.
The elementary characterization and the exact HNF
computation give separate checks of the indices.

The table is supplied as
\texttt{normalized\_order\_indices.tsv}; all nonzero
coordinates of the HNF basis vectors are supplied as
\texttt{normalized\_order\_hnf\_basis.tsv}.  An additional
JSON certificate records the complete matrices and
closure histories.  These files are computation
certificates for this stated model; they are not a
formal proof-assistant certificate or a comparison
theorem for an unspecified automorphic realization.



\section{The cancelled openness remark in cote 34}
\label{sec:q03}

The literal assertion in cote 34, scan 064, Remark 1.17, that
\(\overline{[\pi_0,H]}\) is open in \(\pi'\), is false.
There are two separate issues: the commutators lie in the geometric
kernel \(H\), and a quotient \(\pi_0\) acts on a nonabelian kernel only
through outer automorphisms unless an action or a section is specified.
We prove the natural precise corrections, including the soluble
nonabelian statement.  All openness conclusions below concern \(H\).
The exact intended wording of the crossed-out sentence remains a
question about the original image; a forced mathematical correction is
not a verification of handwriting.

\subsection{The theorem already available in the source}

Let \(k\) be finitely generated over its prime field, let \(X/k\) be a
normal geometrically connected scheme locally of finite type, and let
\(\ell\ne\operatorname{char}(k)\).  For an open subgroup
\(\pi\subseteq\pi_1(X)\), write
\[
1\longrightarrow I\longrightarrow\pi\longrightarrow\Gamma
\longrightarrow1,
\qquad
I=\pi\cap\pi_1(X_{k^{\mathrm{sep}}}).
\]
Thus \(\Gamma\), denoted \(\pi_0\) in the source, is the arithmetic
image.  For a profinite group \(J\), let \(J(\ell)\) denote its maximal
\emph{abelian} pro-\(\ell\) quotient.  Commutator subgroups and
coinvariant relations are always closed topologically.

The background theorem used in this section is precisely cote 34,
scan 053, Theorem 1.1:
\[
 I(\ell)\text{ is a finitely generated }\mathbf Z_\ell\text{-module},
 \qquad I(\ell)_\Gamma\text{ is finite}.
 \tag{Q3-B01}
\]
Its quantifier over \emph{every open subgroup} \(\pi\) is essential.
We carry this established theorem as background, rather than claiming
to have reconstructed the omitted proof pages 065--073.  The supplied
scans 053--063 were read before Remark 1.17; scan 074 was also checked.

Let \(N_0\subseteq I\) be closed and normal in \(\pi\), and suppose
\[
H=I/N_0
\]
is pro-\(\ell\).  Put \(P=\pi/N_0\).  The extension under discussion is
\[
1\longrightarrow H\longrightarrow P\longrightarrow\Gamma
\longrightarrow1. \tag{Q3-E}
\]
Here \(P\) is the \(\pi'\) of the cancelled remark.  Distinguishing this
symbol from the abelian quotient extension used earlier in Theorem 1.1
avoids an ambiguity in the source.

\begin{lemma}[Q3-03: relative finite abelianization]
\label{q3:relative}
For every open subgroup \(P_1\subseteq P\), let \(J=P_1\cap H\).
Then the coinvariant group
\[
(J^{\mathrm{ab}})_{P_1}
\]
for the conjugation action is finite.
\end{lemma}

\begin{proof}
Let \(\pi_1\) be the inverse image of \(P_1\) in \(\pi\).  It is open
in \(\pi_1(X)\).  Its geometric kernel
\(I_1=\pi_1\cap I\) maps onto \(J\): a lift in \(\pi_1\) of an
element of \(J\) has trivial arithmetic image and so belongs to
\(I_1\).  Since \(J\) is pro-\(\ell\), the composite
\[
I_1\twoheadrightarrow J\twoheadrightarrow J^{\mathrm{ab}}
\]
factors through \(I_1(\ell)\).  It is equivariant for \(\pi_1\).
On \(I_1(\ell)\) the conjugation action factors through the arithmetic
image \(\Gamma_1\) of \(\pi_1\).  Consequently there is a surjection
\[
I_1(\ell)_{\Gamma_1}\twoheadrightarrow
(J^{\mathrm{ab}})_{P_1}.
\]
The source is finite by Theorem 1.1 applied to the open subgroup
\(\pi_1\).  This proves the assertion.  No section or solvability
assumption was used.
\end{proof}

\subsection{Literal refutation and the canonical correction}

\begin{proposition}[Q3-01: the stated ambient group is impossible]
\label{q3:literal}
Even when a continuous section of \((\mathrm{Q3\text{-}E})\) is given
and \(H\) is soluble of finite derived length,
\(\overline{[\Gamma,H]}\) need not be open in \(P\).
\end{proposition}

\begin{proof}
Take \(k=\mathbf F_q\), \(X=\operatorname{Spec}k\),
\(\pi=P=\Gamma=G_k=\widehat{\mathbf Z}\), and \(I=H=1\).
The trivial geometric quotient is a soluble pro-\(\ell\) group.
Its commutator subgroup with \(\Gamma\) is trivial, hence is not
open in the infinite profinite group \(P\).  These data meet every
displayed hypothesis.

There is also a nontrivial geometric example.  Take
\(X=\mathbf G_m/\mathbf F_2\), \(\ell=3\), and the rational point
\(1\).  The Kummer tower gives
\[
H=\mathbf Z_3(1),
\]
and the rational point gives a section.  Arithmetic Frobenius acts
on \(H\) by multiplication by \(2\).  Therefore its commutators
already contain \((2-1)H=H\), so
\(\overline{[\Gamma,H]}=H\).  But
\(P/H=\widehat{\mathbf Z}\) is infinite, and \(H\) is not open in
\(P\).
\end{proof}

For a general extension \((\mathrm{Q3\text{-}E})\), conjugation gives
a canonical homomorphism
\[
\Gamma\longrightarrow\operatorname{Out}(H).
\]
Choosing different lifts in \(P\) changes the action on \(H\) by
inner automorphisms.  Thus \([\Gamma,H]\), construed as the
commutator subgroup for an actual action, needs additional data.
The subgroup \([P,H]\) does not.

\begin{proposition}[Q3-04: canonical openness]
\label{q3:canonical}
For every pro-\(\ell\) quotient \(H\) in
\((\mathrm{Q3\text{-}E})\), the subgroup
\[
C=\overline{[P,H]}
\]
is open in \(H\).  No solvability assumption is needed.
\end{proposition}

\begin{proof}
Since \(H\) is normal in \(P\), the closed subgroup \(C\) is normal
in \(P\).  It contains \(\overline{[H,H]}\), so \(H/C\) is abelian.
By definition, \(P\) acts trivially on this quotient.  Hence the
quotient map \(I\to H/C\) factors through
\[
I(\ell)_\Gamma.
\]
This last group is finite by Theorem 1.1, and therefore \(H/C\) is
finite.  A closed subgroup of finite index in a profinite group is
open.
\end{proof}

Equivalently,
\[
H/\overline{[P,H]}=(H^{\mathrm{ab}})_\Gamma.
\]
This correction is exactly the quotient already controlled by
Theorem 1.1.  It does not require a new theorem about purity in
arbitrary dimension.

\subsection{The action-dependent soluble statement}

Suppose now that a continuous section \(s:\Gamma\to P\) is supplied.
Write
\[
\alpha_\gamma(h)=s(\gamma)h s(\gamma)^{-1},\qquad
\delta_\gamma(h)=h^{-1}\alpha_\gamma(h),
\]
and define
\[
N=\overline{\left\langle
 \delta_\gamma(h):h\in H,\ \gamma\in\Gamma
\right\rangle}.
\]
The word soluble in the next result means that the \emph{closed}
derived series terminates after finitely many steps.

\begin{proposition}[Q3-05: normality and the FAb quotient]
\label{q3:fab}
The subgroup \(N\) is normal in \(P\).  For \(Q=H/N\), every open
subgroup \(V\subseteq Q\) has finite abelianization.
\end{proposition}

\begin{proof}
The normality assertion must be checked before forming the
nonabelian quotient.  For \(h,k\in H\), direct multiplication gives
\[
\delta_\gamma(hk)
=k^{-1}\delta_\gamma(h)k\,\delta_\gamma(k).
\]
Thus
\[
k^{-1}\delta_\gamma(h)k
=\delta_\gamma(hk)\delta_\gamma(k)^{-1}.
\]
The subgroup generated by the \(\delta_\gamma(h)\), and therefore
its closure \(N\), is normal in \(H\).  For \(\eta\in\Gamma\),
\[
\alpha_\eta\bigl(\delta_\gamma(h)\bigr)
=\delta_{\eta\gamma\eta^{-1}}\bigl(\alpha_\eta(h)\bigr).
\]
It follows that \(N\) is stable under \(s(\Gamma)\), and hence is
normal in \(P=H\rtimes s(\Gamma)\).

The action of \(\Gamma\) on \(Q=H/N\) is trivial.  The given section
therefore identifies the quotient extension with the direct product
\[
P/N\simeq Q\times\Gamma.
\]
Let \(V\) be an arbitrary open subgroup of \(Q\), not necessarily
normal, and let \(P_V\) be the inverse image of \(V\times\Gamma\)
in \(P\).  The subgroup \(P_V\) is open, and
\[
J_V=P_V\cap H\twoheadrightarrow V.
\]
Conjugation by \(V\times\Gamma\) acts trivially on
\(V^{\mathrm{ab}}\): the first factor acts by inner automorphisms,
and the second acts trivially already on \(V\).  Therefore
\(V^{\mathrm{ab}}\) is a quotient of
\[
(J_V^{\mathrm{ab}})_{P_V}.
\]
Lemma~\ref{q3:relative} makes this group finite.  Thus every open
subgroup \(V\) of \(Q\) has finite abelianization, as asserted.
\end{proof}

A profinite group with the conclusion of
Proposition~\ref{q3:fab} is called \emph{FAb}.  Finiteness of
\(Q^{\mathrm{ab}}\) alone would be insufficient for the next
argument; the quantifier over all open subgroups is what permits
iteration.

\begin{lemma}[Q3-06: soluble FAb groups]
\label{q3:soluble}
A profinite FAb group of finite derived length is finite.
Consequently, if \(H\) in Proposition~\ref{q3:fab} has finite
derived length, then \(N=\overline{[s(\Gamma),H]}\) is open in \(H\).
\end{lemma}

\begin{proof}
For an FAb group \(Q\), let
\[
D^0Q=Q,\qquad
D^{r+1}Q=\overline{[D^rQ,D^rQ]}.
\]
The quotient \(Q^{\mathrm{ab}}\) is finite, so \(D^1Q\) is open.
If \(D^rQ\) is open in \(Q\), FAb says that its abelianization is
finite.  Hence \(D^{r+1}Q\) is open in \(D^rQ\), and therefore in
\(Q\).  Induction makes every \(D^rQ\) open.  Finite derived length
gives \(D^mQ=1\) for some finite \(m\).  Thus \(Q\) is compact and
discrete, and hence finite.

In Proposition~\ref{q3:fab}, the quotient \(Q=H/N\) has no greater
derived length than \(H\).  It is FAb by that proposition, so it is
finite.  This is exactly the openness of \(N\) in \(H\).
\end{proof}

\begin{remark}
Replacing finite derived length by prosolvability would erase a
necessary hypothesis.  Every pro-\(\ell\) group is already
prosolvable, whereas the proof above requires a finite terminal
derived index.  No such extension is asserted here.
\end{remark}

\subsection{An intrinsic variant without a section}

\begin{proposition}[Q3-07: trivial outer arithmetic action]
\label{q3:outer}
Suppose \(H\) in \((\mathrm{Q3\text{-}E})\) has finite derived
length and the canonical map
\(\Gamma\to\operatorname{Out}(H)\) is trivial.  Then \(H\) is finite.
\end{proposition}

\begin{proof}
Let \(C=C_P(H)\), the closed centralizer of \(H\) in \(P\).
Triviality of the outer action means that every element of \(P\)
acts on \(H\) by conjugation by an element of \(H\).  Thus
\[
P=HC.
\]
Let \(D_i\) be the closed derived series of \(H\).  We prove
inductively that \(D_i\) is open in \(H\).  The claim holds for
\(D_0=H\).  Suppose it holds for \(D_i\).  The group \(D_i\) is
characteristic in \(H\), hence normal in \(P\).  The product
\[
R_i=D_i C
\]
is a closed subgroup of \(P\).  Since \(P=HC\),
\[
[P:R_i]\le[H:D_i]<\infty,
\]
so \(R_i\) is open.

Choose an open normal subgroup \(U_i\subseteq P\) such that
\(U_i\cap H\subseteq D_i\).  Such a subgroup exists because \(D_i\)
is an open neighbourhood of \(1\) in the subspace \(H\), and open
normal subgroups form a neighbourhood basis at \(1\) in \(P\).
Since \(D_i\) is normal in \(P\), \(D_iU_i\) is open and
\[
(D_iU_i)\cap H=D_i.
\]
Put
\[
P_i=R_i\cap D_iU_i.
\]
It is open in \(P\), contains \(D_i\), and satisfies
\(P_i\cap H=D_i\).  Every element of \(P_i\subseteq D_iC\) acts on
\(D_i\) by conjugation by an element of \(D_i\), because \(C\)
centralizes \(H\).  Its action on \(D_i^{\mathrm{ab}}\) is
therefore trivial.  Lemma~\ref{q3:relative}, applied to \(P_i\),
shows that \(D_i^{\mathrm{ab}}\) is finite.  Hence \(D_{i+1}\)
is open in \(D_i\), and therefore open in \(H\).

The induction reaches \(D_m=1\) after finitely many steps.  Thus
\(1\) is open in \(H\), proving that \(H\) is finite.
\end{proof}

This is an optional intrinsic result about the extension; it does
not identify the intended wording of Remark 1.17.

\subsection{The exact weight input and the role of Weil I}

\begin{lemma}[Q3-02: a Frobenius without eigenvalue \(1\)]
\label{q3:weights}
Let \(M\) be a finitely generated \(\mathbf Z_\ell\)-module with
continuous action by a profinite group \(G\).  If some \(g\in G\)
has no eigenvalue \(1\) on \(M\otimes_{\mathbf Z_\ell}\mathbf Q_\ell\),
then \(M_G\) is finite.
\end{lemma}

\begin{proof}
The module \(M/(g-1)M\) is finitely generated over
\(\mathbf Z_\ell\).  Its tensor product with \(\mathbf Q_\ell\)
vanishes because \(g-1\) is invertible after tensoring.  Thus it is
a finitely generated torsion \(\mathbf Z_\ell\)-module and is
finite.  It surjects onto \(M_G\), so \(M_G\) is finite.
\end{proof}

If \(g\) is a finite-field Frobenius with pure eigenvalues of a
nonzero weight, the eigenvalue \(1\) is excluded.  Replacing
geometric Frobenius by arithmetic Frobenius inverts eigenvalues
and reverses weights, without affecting that exclusion.
Semisimplicity is not needed.

For a smooth projective curve \(C/\mathbf F_q\), Weil I,
\cite[(1.6)--(1.7)]{WeilI}, gives weight \(1\) on
\(H^1(C_{\overline{\mathbf F}_q},\mathbf Q_\ell)\), hence weight
\(-1\) on its dual.  The identification
\[
H^1(C_{\overline{\mathbf F}_q},\mathbf Q_\ell)
=\operatorname{Hom}_{\mathrm{cont}}
 \bigl(\pi_1(C_{\overline{\mathbf F}_q}),\mathbf Q_\ell\bigr)
\]
follows from the classification of finite constant-group torsors
and passage to the inverse limit.  This supplies the
no-eigenvalue-\(1\) input for Lemma~\ref{q3:weights}.

The spreading and specialization mechanism is explicit in cote
34, scan 074, formulas (2.5.4)--(2.5.6): extend the smooth
projective curve over a finitely generated normal ring, compare
the lisse \(H^1\)-fibres, and specialize at a closed point with
finite residue field.  One Frobenius without eigenvalue \(1\)
suffices.  Scan 054 announces the reduction to curves, and scan
053 states the resulting Theorem 1.1.  This is why the curve
case contained in Weil I is sufficient for the established input
used here.  We are not asserting that smooth projective purity
alone, without those geometric reductions, proves an arbitrary
normal-scheme fundamental-group statement.

An independently accessible published account of the same
mechanism is \cite{KatzLang}: Theorem 1 and its proof in
Section II, followed by Theorems \(1(\mathrm{bis})\) and
\(1(\mathrm{ter})\).  The proof passes through Artin good
neighbourhoods and generalized Jacobians to abelian varieties
and tori over finite fields; its exact cited locators for the
neighbourhood reduction are SGA 4 XI 3.1 and 3.3.  This confirms
the curve/Jacobian/spreading mechanism, but is not substituted
without further argument for the full normal-scheme formulation
of the supplied Theorem 1.1.

For comparison, \cite[(3.3.5), p.~206]{WeilII} supplies a broader
direct weight estimate: on a smooth variety over a finite field,
\(H^1\) has weights at least \(1\), so its dual has weights at
most \(-1\).  Lemma~\ref{q3:weights} then applies.  This Weil II
route is unnecessary for
Propositions~\ref{q3:canonical}, \ref{q3:fab}, and
\ref{q3:outer}; it records accurately the distinction between
smooth projective purity and a nonproper weight estimate.

The dependency chain for the corrected statements is therefore
\[
\begin{gathered}
\text{curve RH, contained in Weil I}
+\text{ the geometric reductions used in Theorem 1.1}\\
\Longrightarrow
\text{Theorem 1.1 for every open }\pi
\Longrightarrow \text{Q3-03}\\
\Longrightarrow
\text{Q3-04, or Q3-05--Q3-06, or Q3-07}.
\end{gathered}
\]
None of these implications uses a higher-dimensional purity
theorem, a Hodge standard conjecture, or the assumed arithmetic
RH/Dirichlet-GRH input.  The literal assertion in
Proposition~\ref{q3:literal} remains false under all such
additional hypotheses.

\subsection{The locally finite type qualification}

\begin{lemma}[Q3-08: reduction to a nonempty affine open]
\label{q3:affine}
Under the exact-sequence framework stated in Theorem 1.1, its
finite-generation and finite-coinvariant conclusions for normal
geometrically connected finite-type schemes imply the same
conclusions for normal geometrically connected schemes locally
of finite type.
\end{lemma}

\begin{proof}
A connected normal locally noetherian scheme is integral:
normality makes distinct irreducible components disjoint, and
local noetherianity makes each component locally open.
Connectedness leaves a single component.  Choose a nonempty
affine open \(U\subseteq X\).  It is finite type over \(k\).

If \(X'\to X\) is a connected finite étale cover, then \(X'\)
is normal and connected, hence integral.  Its inverse image of
\(U\) is a nonempty open subset of an integral scheme and is
therefore connected.  The connected-cover criterion
\cite[V, Proposition 6.9]{SGA1} gives a surjection
\[
\pi_1(U)\twoheadrightarrow\pi_1(X).
\]
Normality is preserved by separable algebraic field extension.
Applying the same argument over \(k^{\mathrm{sep}}\), where the
assumed geometric connectedness applies, gives
\[
\pi_1(U_{k^{\mathrm{sep}}})
\twoheadrightarrow\pi_1(X_{k^{\mathrm{sep}}}).
\]

For an open subgroup \(\pi\subseteq\pi_1(X)\), let \(\pi_U\)
be its inverse image in \(\pi_1(U)\).  Its geometric kernel
\(I_U\) maps onto \(I\), and both arithmetic images are the same
group \(\Gamma\).  Consequently
\[
I_U(\ell)\twoheadrightarrow I(\ell),\qquad
I_U(\ell)_\Gamma\twoheadrightarrow I(\ell)_\Gamma.
\]
The finite-type conclusion for \(U\) therefore transfers finite
generation and finiteness of coinvariants to \(X\).
\end{proof}

The published exact-sequence locator is
\cite[IX, Theorem 6.1]{SGA1}, originally p.~253
(p.~195 in the re-edition); that statement includes a
quasi-compactness hypothesis.  The reference to Exposé VIII in
scan 053 does not match this published locator and should be
flagged as a citation candidate, possibly reflecting an earlier
numbering.  Lemma~\ref{q3:affine} does not silently apply the
quasi-compact theorem to a non-quasi-compact scheme: it uses the
exact-sequence framework already supplied in Theorem 1.1 and
proves the surjections needed to transfer its conclusions.
The underlying Galois-category equivalence is
\cite[V, Theorem 4.1]{SGA1}.

\subsection{Disposition of the question}

The literal row is refuted by Q3-01.  The canonical correction
\(\overline{[P,H]}\) open in \(H\) is Q3-04.  With a specified
continuous section and finite derived length,
\(\overline{[s(\Gamma),H]}\) open in \(H\) is Q3-05--Q3-06.
The intrinsic trivial-outer-action variant is Q3-07.
These statements are settled within the established
fundamental-group and curve-RH framework.  The remaining textual
question is which correction, if any, the author intended in the
cancelled Remark 1.17.



\part{Unestablished stronger claims and source boundaries}
\section{What the checked results do not establish}

This part contains no finding presented as a theorem.  It records the
precise propositions that the task's suggestions could otherwise be read
as asserting, and the limits of the supplied evidence.  A counterexample
to a formal inference is not a counterexample to a geometric conjecture.

\begin{longtable}{@{}p{0.12\textwidth}p{0.35\textwidth}p{0.46\textwidth}@{}}
\toprule
ID & Unestablished assertion & Exact boundary or required additional datum\\
\midrule
\endfirsthead
\toprule ID & Unestablished assertion & Exact boundary or required additional datum\\
\midrule\endhead
U-01 & General Hodge standard positivity, or algebraic Lefschetz correspondences, from sign-only purity. & HS-05 and SG-01--03 refute the proposed formal implication. A construction of the specified positive form and its algebraic adjoint remains additional geometry. BRIDGES/Q05 is unchanged.\\
U-02 & Every algebraic class on an arbitrary CM abelian variety is a Lefschetz class. & SG-04 proves this only for the specified CM elliptic powers and their good reductions. Milne's conditional general theorem retains its actual Hodge-conjecture hypothesis.\\
U-03 & A higher arithmetic Hodge-index theorem, or a corrected general cote-44 height identity, from steps 5 and 13. & AR-01 refutes the displayed raw identity; AR-02--03 refute the proposed sign inference. A corrected theorem must specify admissible cycle extensions, vertical corrections, metrics when arithmetic, and comparison with the canonical self-duality.\\
U-04 & A normalized, lattice-preserving identification of the intended automorphic lift chain with a Lefschetz realization. & SL-05 and SL-07 rule out the automatic lattice-free and untwisted identifications. A new comparison would require a specified realization, Tate normalization, integral intertwiner, and equality of the actual operators and lattices. No such comparison is asserted.\\
U-05 & Exact valuations of a particular nonzero iterated theta period determined solely by the lift count. & SL-06 gives the actual conditional local bounds and explicit numerical contributions. Its inputs include highest weights, primes, denominator bounds and entry estimates. Wilson computes an entirely specified exterior map; the example does not select a nonzero period.\\
U-06 & The original full SA06 C.1f--q proofs, the degree-14 Lawson nonvanishing in E.3u, or the original handwriting in the struck remarks. & Only the supplied intake statements and transcriptions were available. They are not independently upgraded. The present purity assertions do not depend on those absent proofs; the proposed textual repairs require the images.\\
U-07 & A closed Lean certificate for the reconstructed generic-path proof. & The archive was inspected at declaration and selected proof-body level, not built. Source definitions and implication theorems are distinguished from supplied premises. All theorem proofs used here are mathematical proofs with their explicit foundational inputs.\\
\bottomrule
\end{longtable}

The distinction in U-04--05 is part of the answer to the comparison
question: the stated unqualified identification is false, and the two
actual numerical statements have been computed with their separate
parameters.  An additional normalized integral comparison has not been
invented, assumed, or disproved merely because it is not supplied.

\section{Dependencies, verification, and programme consequences}

The main implication chain is
\[
 \begin{gathered}
 \text{Artin vanishing + generic weak Lefschetz + duality}\\
 \Downarrow\\
 \text{SG-08: an effective Gysin kernel and a tensor-power summand}\\
 \text{Noether normalization + HS-03 (Parseval) + HS-01--02}\\
 \Downarrow\\
 \text{SG-09: affine middle upper bound}\\
 \Downarrow\\
 \text{SG-11: compact-support upper bounds}\\
 \Downarrow\\
 \text{SG-10: smooth proper purity by duality.}
 \end{gathered}
\]
There is no arrow here from a trace multiplicity to a positive
intersection form.  Cote 16's evanescent machinery helps locate the
effective kernel; its index and cycle-conjecture assertions are separate.
The only use of the stipulated global RH in these findings is the
change of variable for the explicitly identified Tate $L$-factor in
SG-02.  No general positive-weight number-field $L$-function is converted
into a Dirichlet $L$-function by this argument.

The original BRIDGES/Q03 row is replaced by a literal refutation and
proved precise corrections (Q3-01, Q3-04--07).  BRIDGES/Q04 remains the
known weight-to-coinvariants consequence.  BRIDGES/Q05 remains open in
its general form.  The present SA06 C.1 conclusions are realization
statements: they leave all independent cycle-class-kernel issues in
their proper scope.  H8's exterior-algebra content is corrected and
computed in SL-01--04 and SL-08.  H9's displayed uncorrected higher-cycle
identity is refuted by AR-01.  No statement becomes conditional on
Masullo here, since no specific Masullo theorem is an input.

Every complete proof in Part I received a full reading after its
construction; the principal sign-only theorem and the main counterexamples
also received independent reconstruction by another agent.  The exact
events, including corrected drafts, appear in \texttt{CHECK\_EVENTS.tsv}.
All agents used the inherited GPT-6 Astra Pro model; the runtime did not
expose a separate effort setting, and the log does not manufacture one.
The integer matrix checks are reproducible with the included Python
script and tables.  They are computations, not Lean certificates.
Literature statements inspected but not re-proved are logged with that
weaker event type.  The dated ledger records actual gate changes;
preliminary notes are preserved as working history and explicitly
superseded where the completed geometric proof closed a former gap.

\begin{thebibliography}{99}
\bibitem{Alpoge}
L.~Alp\"oge, \emph{Remark}, supplied four-page note dated 20 November 2026.
Step 5: p.~3, footnote 3 (timeline steps 3 and 5); step 10: divisibility
of the iterated lift; step 13: slope inequality.
\url{https://alpo.ge/remark.pdf}. Supplied wrappers
\texttt{ALPOGE\_Main.lean}, \texttt{ALPOGE\_Statements.lean}; retrieved
source/declaration locators are given in the body and evidence notes.

\bibitem{Archive}
A.~Grothendieck, supplied archive transcriptions: cote 16, scans 014--018,
034--040, 049--050; cote 44, scans 001--009; cote 34, scans 053--064 and
074. The exact file inventory and hashes are in \texttt{SOURCE\_MANIFEST.tsv}.
Original scan images were not included.

\bibitem{SGA1}
A.~Grothendieck et al., \emph{Rev\^etements \'etales et groupe fondamental}
(SGA 1), LNM 224, Springer, 1971; corrected re-edition,
\url{https://arxiv.org/pdf/math/0206203}.
Expos\'e V, Th.~4.1 and Prop.~6.9; Expos\'e IX, Th.~6.1.

\bibitem{SGA4}
M.~Artin, A.~Grothendieck, J.-L.~Verdier et al., \emph{Th\'eorie des topos
et cohomologie \'etale des sch\'emas} (SGA 4), LNM 269, 270, 305.
Expos\'e XI, 3.1, 3.3; XIV, Th.~3.1; XVII, 5.4.3; XVIII (duality).
The uses and source-check depth are recorded in the check table.

\bibitem{SGA4half}
P.~Deligne et al., \emph{Cohomologie \'etale} (SGA $4\frac12$), LNM 569,
Springer, 1977, \emph{Th\'eor\`emes de finitude}, 1.9 and 4.3.

\bibitem{BBD}
A.~Beilinson, J.~Bernstein, P.~Deligne, \emph{Faisceaux pervers},
Ast\'erisque 100 (1982), 4.1.1--4.1.3 and 4.2.8.
\url{https://publications.ias.edu/sites/default/files/Faisceaux%20pervers.pdf}.

\bibitem{WeilI}
P.~Deligne, \emph{La conjecture de Weil. I}, Publ. Math. IHES 43 (1974),
273--307. (1.5.1), (1.5.4), Th.~1.6, (1.7), (2.3), (2.14),
\S\S4.1, 5.9.
\url{https://www.numdam.org/article/PMIHES_1974__43__273_0.pdf}.

\bibitem{WeilII}
P.~Deligne, \emph{La conjecture de Weil. II}, Publ. Math. IHES 52 (1980),
137--252, Th.~3.3.1 and Cor.~3.3.5.
\url{https://publications.ias.edu/sites/default/files/Number40.pdf}.

\bibitem{KatzACT}
N.~M.~Katz, \emph{Affine Cohomological Transforms, Perversity, and
Monodromy}, J. Amer. Math. Soc. 6 (1993), 149--222.
Th.~3.4, Cor.~3.4.1, Cor.~3.4.3 and Lemma 3.5; Appendix A.3--A.5.
\url{https://web.math.princeton.edu/~nmk/old/ACT.pdf}.

\bibitem{StacksNN}
The Stacks Project, \emph{Noether normalization}, Lemma 10.115.4,
Tag 00OY, \url{https://stacks.math.columbia.edu/tag/00OY}.

\bibitem{KatzLang}
N.~M.~Katz and S.~Lang, \emph{Finiteness theorems in geometric classfield
theory}, Enseign. Math. 27 (1981), 285--319, Th.~1, proof in \S II,
Ths.~1(bis), 1(ter).
\url{https://web.math.princeton.edu/~nmk/old/AbsFinThm.pdf}.

\bibitem{MilnePolarizations}
J.~S.~Milne, \emph{Polarizations and Grothendieck's Standard Conjectures},
Ann. of Math. 155 (2002), 599--610. Published-version preprint,
Prop.~3.2, Th.~3.3, Remarks 3.7--3.8.
\url{https://www.jmilne.org/math/articles/2002aP.pdf}.

\bibitem{MilneAV}
J.~S.~Milne, \emph{Abelian Varieties}, version 1.10 (2008).
Th.~2.3; Th.~11.1(b); Prop.~13.1; Th.~13.2; Th.~14.1 and
Lemmas 14.2--14.3.
\url{https://www.jmilne.org/math/CourseNotes/AV110.pdf}.

\bibitem{SutherlandEC}
A.~V.~Sutherland, \emph{18.783 Elliptic Curves}, MIT OpenCourseWare,
Fall 2025, Lecture 13, Ths.~13.3, 13.18 and Cor.~13.20.
\url{https://ocw.mit.edu/courses/18-783-elliptic-curves-fall-2025/mit18_783_f25_lec_full.pdf}.

\bibitem{MilneJacobian}
J.~S.~Milne, \emph{Jacobian Varieties}, Chapter VII of
\emph{Arithmetic Geometry}, Springer, 1986; author's corrected version
of 12 June 2021, Theorem 6.6 and Summary 6.11. \url{https://www.jmilne.org/math/xnotes/JVs.pdf}.

\bibitem{Bost}
J.-B.~Bost, \emph{Algebraic leaves of algebraic foliations over number
fields}, Publ. Math. IHES 93 (2001), 161--221,
\S4.1.4, Props.~4.5--4.6, equations (4.11), (4.14).
\url{https://www.numdam.org/article/PMIHES_2001__93__161_0.pdf}.

\bibitem{Fulton}
W.~Fulton, \emph{Intersection Theory}, second edition, Springer, 1998,
Th.~3.3(b), projective bundle formula.

\bibitem{GilletSoule}
H.~Gillet and C.~Soul\'e, \emph{Arithmetic intersection theory},
Publ. Math. IHES 72 (1990), 93--174, \S4.2.11, \S4.3,
Th.~4.3.2.
\url{https://www.numdam.org/article/PMIHES_1990__72__93_0.pdf}.

\bibitem{Moriwaki}
A.~Moriwaki, \emph{Hodge index theorem for arithmetic cycles of
codimension one}, Theorems A and B, arXiv:alg-geom/9403011v4.
\url{https://arxiv.org/pdf/alg-geom/9403011}.

\bibitem{Kunnemann}
K.~K\"unnemann, \emph{Some remarks on the arithmetic Hodge index
conjecture}, Compositio Math. 99 (1995), 109--128,
Th.~4.1 and \S5.
\url{https://www.numdam.org/article/CM_1995__99_2_109_0.pdf}.

\bibitem{Shioda}
T.~Shioda, \emph{On the Mordell--Weil lattices}, Comment. Math. Univ.
St. Pauli 39 (1990), 211--240, Lemma 8.1, Ths.~8.4, 8.6,
equations (8.8), (8.11)--(8.14).
\url{https://rikkyo.repo.nii.ac.jp/record/10027/files/AA00610867_39-02_05.pdf}.

\bibitem{Wilson}
R.~M.~Wilson, \emph{A diagonal form for the incidence matrices of
$t$-subsets vs.\ $k$-subsets}, European J. Combin. 11 (1990),
609--615. Used through the verified Corollary 3 of the next reference.

\bibitem{Ghorbani}
E.~Ghorbani, G.~B.~Khosrovshahi, Ch.~Maysoori, M.~Mohammad-Noori,
\emph{Inclusion Matrices and Chains}, arXiv:0709.3144v1,
Corollary 3, p.~10; J. Combin. Theory Ser. A 115 (2008), 878--887.
\url{https://arxiv.org/pdf/0709.3144}.

\bibitem{Mok}
C.~P.~Mok, \emph{Endoscopic Classification of Representations of
Quasi-Split Unitary Groups}, arXiv:1206.0882v5,
\S2.2, equations (2.2.9)--(2.2.10), and \S2.3, equation (2.3.1).
\url{https://arxiv.org/pdf/1206.0882}.

\bibitem{MilneEtale}
J.~S.~Milne, \emph{Lectures on \'Etale Cohomology}, version 2.21,
Th.~23.3, \S27 and Prop.~27.3.
\url{https://www.jmilne.org/math/CourseNotes/LEC.pdf}.
\end{thebibliography}
\end{document}
